% 褐矮星（综述）
% license CCBYSA3
% type Wiki

本文根据 CC-BY-SA 协议转载翻译自维基百科\href{https://en.wikipedia.org/wiki/Brown_dwarf}{相关文章}。

\begin{figure}[ht]
\centering
\includegraphics[width=6cm]{./figures/6432b0e71e12ceea.png}
\caption{T 型褐矮星的艺术概念图} \label{fig_Brdw_1}
\end{figure}
\begin{figure}[ht]
\centering
\includegraphics[width=6cm]{./figures/dc30c14f9295e490.png}
\caption{对比：大多数褐矮星的体积仅比木星略大（约大15–20\%）$^{[1]}$，但由于其密度更高，质量仍可达到木星的 80 倍。图像按比例绘制，其中木星的半径约为地球的 11 倍，而太阳的半径约为木星的 10 倍。} \label{fig_Brdw_2}
\end{figure}
褐矮星是亚恒星天体，其质量大于最大的气态巨行星，但小于质量最小的主序星。它们的质量大约为木星的 13–80 倍（($M_J$)）$^{[2]}$——不足以在其核心中维持将氢聚变为氦的核反应，但又足够巨大，可以通过氘（$^2\mathrm{H}$）的聚变释放一定的光和热。氘是氢的一种同位素，含有一个中子和一个质子，其聚变所需温度更低。质量最大的褐矮星（$(>65,M_J)$）还能够进行锂（$(^7\mathrm{Li})$）的聚变。

天文学家按照光谱型对自发光天体进行分类，这一分类与表面温度密切相关。褐矮星分布在 M 型（2100–3500 K）、L 型（1300–2100 K）、T 型（600–1300 K）和 Y 型（<600 K） 光谱类别中 $^{[3][4]}$。由于褐矮星不会发生稳定的氢聚变，它们会随着时间推移逐渐冷却，并在演化过程中依次进入更晚的光谱型。

“褐矮星”中的“褐色”原意是指介于红色与黑色之间的一种颜色 $^{[5]}$。以肉眼观察，大多数褐矮星呈现为品红色，而其他褐矮星则会因温度不同而显示出不同的颜色 $^{[3][6]}$。褐矮星可能是完全对流的，内部不存在明显的分层结构或随深度变化的化学分异 $^{[7]}$。

尽管早在 20 世纪 60 年代就有人从理论上预言了褐矮星的存在，但直到 1994 年，第一批明确无误的褐矮星才被发现 $^{[8]}$。由于褐矮星的表面温度相对较低，它们在可见光波段并不明亮，而主要在红外波段辐射能量。随着红外探测设备能力的提升，天文学家已经识别出数以千计的褐矮星。目前已知距离最近的褐矮星位于 Luhman 16 系统，这是一个由 L 型和 T 型褐矮星组成的双星系统，距离太阳约 6.5 光年（2.0 秒差距）。在距离上，Luhman 16 是继半人马座 $\alpha$ 星系和巴纳德星之后，距离太阳第三近的恒星系统。
\subsection{历史}
\begin{figure}[ht]
\centering
\includegraphics[width=6cm]{./figures/16619a30e2a82cab.png}
\caption{较小的天体是葛利泽 229B（Gliese 229B），其质量约为木星的 20 至 50 倍，绕恒星葛利泽 229 运行。它位于天兔座，距离地球约 19 光年。} \label{fig_Brdw_3}
\end{figure}
\subsubsection{早期理论研究}
\begin{figure}[ht]
\centering
\includegraphics[width=6cm]{./figures/a5a243a9e6a937ce.png}
\caption{行星、褐矮星、恒星（非按比例）} \label{fig_Brdw_4}
\end{figure}
20 世纪 60 年代，Shiv Kumar 提出了后来被称为褐矮星的天体的存在理论；它们最初被称为黑矮星$^{[9]}$，这一名称用于指在太空中自由漂浮、质量不足以维持氢聚变的暗弱亚恒星天体。然而：\\

(a) “黑矮星”一词此前已被用来指冷却后的白矮星；\\
(b) 红矮星能够进行氢聚变；\\
(c) 这些天体在其生命早期可能在可见波段仍然具有一定亮度。\\

因此，人们提出了其他名称来称呼这些天体，包括 planetar 和 substar。1975 年，Jill Tarter 在加州大学伯克利分校的博士论文中，首次建议使用“褐矮星（brown dwarf）”这一术语，用“褐色”来表示一种“介于红色与黑色之间”的颜色，意指这些矮星看起来昏暗、黯淡、无光泽，尽管它们并非真正呈褐色$^{[5][10][11]}$。

“黑矮星”这一术语至今仍用于指已经冷却到几乎不再发出可探测光辐射的白矮星。然而，即便是最低质量的白矮星，要冷却到这一温度所需的时间也被计算为超过宇宙当前的年龄，因此这类天体预计至今尚未存在$^{[12]}$。

关于最低质量恒星及氢燃烧下限的早期理论认为：第一星族天体若质量小于 $0.07,M_\odot$，或第二星族天体若质量小于 $0.09,M_\odot$，将永远不会经历正常的恒星演化过程，而会成为完全简并的恒星$^{[13]}$。由此产生的褐矮星有时被称为“失败的恒星”$^{[14]}$。首个自洽的氢燃烧最小质量计算结果证实，第一星族天体的下限约为 $0.07-0.08,M_\odot$,$^{[15][16]}$。
\subsubsection{氘聚变}
20 世纪 80 年代末，人们发现质量低至 $0.013,M_\odot$（约 $13.6,M_J$）的天体也可以发生氘聚变，同时褐矮星冷却外层大气中尘埃形成的影响也被认识到，这些发现对上述理论提出了挑战。然而，这类天体极难被发现，因为它们几乎不在可见光波段发射辐射，其最强的辐射位于红外（IR）波段。当时的地基红外探测器精度不足，难以可靠识别褐矮星。

此后，研究者采用多种方法开展了大量搜索工作，包括：对场星的多色成像巡天；对主序矮星和白矮星暗弱伴星的成像巡天；对年轻星团的调查；以及通过径向速度监测寻找近距离伴星。
\subsubsection{GD 165B 与 L 型}

多年来，发现褐矮星的尝试一直未获成功。直到 1988 年，在一次针对白矮星的红外巡天中，人们发现白矮星 GD 165 周围存在一个暗弱伴星。该伴星 GD 165B 的光谱呈现出非常红且难以解释的特征，既不符合低质量红矮星的典型特征，也不同于当时已知的最冷 M 型矮星。这表明，GD 165B 必须被归类为一种比已知 M 型矮星更冷的天体。

在接近十年的时间里，GD 165B 一直是独一无二的例子。直到 1997 年 2MASS（Two Micron All-Sky Survey）启动后，人们才发现了大量在颜色和光谱特征上与其相似的天体。

如今，GD 165B 被认为是后来被称为 L 型矮星的一类天体的原型$^{[17][18]}$。

尽管当时发现如此冷的矮星具有重大意义，但学界一度争论 GD 165B 究竟应被归类为褐矮星，还是仅仅是一颗极低质量恒星，因为在观测上二者极难区分$^{[19][20]}$。

在 GD 165B 被发现之后不久，又陆续有其他褐矮星候选体被报告。然而，其中大多数未能通过检验，因为它们的大气中缺乏锂元素，表明它们实际上是恒星。真正的恒星会在约$10^8$年左右燃尽其锂，而褐矮星（尽管其温度和光度可能与恒星相近）则不会发生这一过程。因此，在年龄超过 $100,\mathrm{Myr}$ 的天体大气中探测到锂元素，几乎可以确定该天体是一颗褐矮星。
\subsubsection{Teide 1 与 M }
首个被确认的 M 型褐矮星于 1994 年由西班牙天体物理学家 Rafael Rebolo（课题组负责人）、María Rosa Zapatero-Osorio 和 Eduardo L. Martín 发现$^{[21]}$。该天体位于昴宿星团这一疏散星团中，被命名为 Teide 1。相关发现论文于 1995 年 5 月投稿至《自然》，并于 1995 年 9 月 14 日正式发表$^{[22][23]}$。在该期《自然》杂志的封面上，醒目标题写道：“褐矮星被发现，官方确认”。

Teide 1 是 IAC 团队于 1994 年 1 月 6 日在特内里费岛泰德天文台，使用 80 厘米口径望远镜（IAC 80）获取的图像中发现的；其光谱则于 1994 年 12 月首次由拉帕尔马岛罗克·德·洛斯·穆查乔斯天文台的 4.2 米威廉·赫歇尔望远镜观测获得。由于 Teide 1 属于年轻的昴宿星团，其距离、化学组成和年龄得以较为准确地确定。利用当时最先进的恒星与亚恒星演化模型，研究团队估计 Teide 1 的质量为$55 \pm 15,M_J$,$^{[24]}$，低于恒星质量下限。此后，Teide 1 成为年轻褐矮星研究中的一个重要参照对象。

从理论上讲，质量低于$65,M_J$的褐矮星在其整个演化过程中都无法通过热核聚变燃烧锂。这一事实构成了锂检验原理的基础之一，用以判断低光度、低表面温度天体是否具有亚恒星性质。

1995 年 11 月，使用 Keck I 望远镜获得的高质量光谱数据显示，Teide 1 仍然保留着形成昴宿星团恒星的原始分子云中的初始锂丰度，从而证明其核心中不存在热核聚变过程。这些观测结果不仅完全确认了 Teide 1 是一颗褐矮星，也验证了光谱锂检验方法的有效性$^{[25]}$。

在相当一段时间内，Teide 1 是通过直接观测所识别的、太阳系外已知质量最小的天体。此后，人们已鉴认出超过 1,800 颗褐矮星$^{[26]}$，其中一些距离地球非常近，例如 Epsilon Indi Ba 与 Bb——一对与一颗类太阳恒星引力束缚、距离太阳约 12 光年的褐矮星系统$^{[27]}$，以及 Luhman 16——一个距离太阳仅 6.5 光年的褐矮星双星系统。
\subsubsection{Gliese 229B 与 T 型}
首个 T 型褐矮星于 1994 年由加州理工学院的天文学家 Shrinivas Kulkarni、中岛敬史（Tadashi Nakajima）、Keith Matthews 和 Rebecca Oppenheimer，以及约翰斯·霍普金斯大学的科学家 Samuel T. Durrance 和 David Golimowski 共同发现$^{[28]}$。该天体于 1995 年被确认是 Gliese 229 的一颗亚恒星伴星，即 Gliese 229b。Gliese 229b 与 Teide 1 一同构成了最早两例清晰的褐矮星观测证据。二者均在 1995 年通过探测到 670.8 nm 的锂谱线得到确认，其中 Gliese 229b 的温度和光度显著低于恒星范围。

其近红外光谱清楚地显示出 2 微米处的甲烷吸收带，这一特征此前仅在巨行星大气以及土星卫星泰坦的大气中被观测到。甲烷吸收并不可能出现在任何主序恒星的温度条件下。正是这一发现促成了比 L 型矮星更冷的新光谱类型——T 型矮星的确立，而 Gliese 229b 则被视为该类型的原型。
\subsection{理论}
\begin{figure}[ht]
\centering
\includegraphics[width=6cm]{./figures/851acc2df0619164.png}
\caption{} \label{fig_Brdw_5}
\end{figure}
恒星形成的标准机制是：一团寒冷的星际气体与尘埃云在自身引力作用下发生坍缩。随着云团收缩，由于 Kelvin–Helmholtz 机制，其温度会升高。在这一过程的早期阶段，收缩中的气体会迅速将大量能量以辐射形式释放出去，从而使坍缩得以继续。最终，中心区域的密度增大到足以俘获辐射的程度。于是，坍缩云团中心的温度和密度会随时间显著上升，收缩速度逐渐减慢，直到核心条件变得足够高温、高密，从而在原恒星核心中触发热核反应。对于一颗典型恒星而言，其核心内由热核聚变反应产生的气体压和辐射压可以支撑恒星抵抗进一步的引力收缩。此时达到流体静力平衡，恒星将在其大部分寿命中以主序星的形式，将氢聚变为氦。

然而，如果原恒星的初始质量$^{[29]}$小于约 (0.08,M_{\odot})$^{[30]}$，则正常的氢热核聚变反应将无法在其核心中点燃。引力收缩并不能非常有效地加热这样一个低质量的原恒星；在核心温度尚未升高到足以触发聚变之前，密度便已上升到电子彼此紧密堆积的程度，从而产生量子意义上的电子简并压。根据褐矮星内部结构模型，核心中典型的密度、温度和压力条件预计如下：
\begin{itemize}
\item 
\[
10,\mathrm{g/cm^3}\ \lesssim\ \rho_c\ \lesssim\ 10^3,\mathrm{g/cm^3},~
\]
\item 
\[
T_c \lesssim 3\times 10^6,\mathrm{K},~
\]
\item 
\[
P_c \sim 10^5,\mathrm{Mbar}.~
\]
\end{itemize}
这意味着，该原恒星在任何时候都不具备维持氢聚变所需的质量或密度条件。正在落入的物质会被电子简并压所阻止，无法达到所需的高密度和高压力。

进一步的引力收缩因此被抑制，最终形成的结果是一颗褐矮星，它仅通过向外辐射自身内部的热能而逐渐冷却。需要注意的是，从原理上讲，褐矮星有可能在不点燃氢聚变的情况下，缓慢吸积物质并使其质量超过氢燃烧下限。这种情况可能通过褐矮星双星系统中的质量转移过程而发生$^{[29]}$。
\subsubsection{高质量褐矮星与低质量恒星}
锂通常存在于褐矮星中，而在低质量恒星中则一般不存在。恒星一旦达到足以发生氢聚变的高温，其内部的锂会很快被消耗殆尽。锂-7 与质子的聚变会生成两个氦-4 原子核，而该反应所需的温度仅略低于氢聚变所需的温度。低质量恒星中的对流作用会使整个恒星内部的锂最终全部被耗尽。因此，在一个候选褐矮星的光谱中探测到锂谱线，是其确为次恒星天体的一个强有力指标。

\textbf{锂检验}

褐矮星可以分为两类：一类具有足够的质量可以发生锂聚变，另一类则不能。这一区分方法被称为锂检验$^{[31]}$。

像太阳这样质量更大的恒星，也可能在其外层保留锂，因为这些外层从未达到足以发生锂聚变的温度，而且其对流层并不会与核心混合；若发生混合，锂将会迅速被消耗。这类较大质量的恒星可以通过其尺寸和光度轻易地区分于褐矮星。

相反地，处在质量上限附近的褐矮星，在其年轻阶段可能足够炽热，从而耗尽自身的锂。质量大于 $65,M_J$的褐矮星，在其约五亿年寿命之前便可以燃烧掉锂$^{[32]}$。

\textbf{大气中的甲烷}

与恒星不同，年龄较老的褐矮星有时会冷却到足够低的温度，使其大气在极长的时间尺度上积累可观测量的甲烷，而甲烷无法在更高温的天体中形成。通过这种方式得到确认的褐矮星包括 Gliese 229B。

\textbf{铁、硅酸盐与硫化物云}

主序星在演化过程中会逐渐冷却，但最终会达到一个能够通过稳定核聚变维持的最小总光度。该光度因恒星而异，但通常至少为太阳光度的 0.01\%$^{[\text{citation needed}]}$。而褐矮星则在其整个寿命中持续冷却并变暗；足够古老的褐矮星将会暗淡到无法被探测的程度。
\begin{figure}[ht]
\centering
\includegraphics[width=8cm]{./figures/963216b7bed1d372.png}
\caption{早期 T 型褐矮星 SIMP J0136+09 和 2MASS J2139+02（左侧两个面板）的云模型，以及晚期 T 型褐矮星 2M0050–3322 的云模型。} \label{fig_Brdw_6}
\end{figure} 
云被用来解释晚期 L 型褐矮星中铁氢化物（FeH）光谱线的减弱。铁云会使上层大气中的 FeH 被耗尽，同时云层遮挡了对仍然含有 FeH 的更深层大气的观测视线。随后，在中期到晚期 T 型褐矮星更低的温度条件下，这种化学物质再次增强，被解释为云层受到扰动，使望远镜能够观测到仍然含有 FeH 的更深大气层$^{[33]}$。年轻的 L/T 型褐矮星（L2–T4）表现出很强的光变，这可能可以用云层、热点、磁驱动的极光或热化学不稳定性来解释$^{[34]}$。这些褐矮星的云被解释为厚度变化的铁云，或是下方一层厚铁云与上方一层硅酸盐云的组合。上层硅酸盐云可能由石英（quartz）、顽火辉石（enstatite）、刚玉（corundum）和／或橄榄石（forsterite）组成$^{[35][36]}$。然而，目前尚不清楚硅酸盐云是否对所有年轻天体都是必需的$^{[37]}$。硅酸盐吸收可以在 8–12 μm 的中红外波段被直接观测到。利用 Spitzer 的 IRS 进行的观测表明，硅酸盐吸收在 L2–L8 型褐矮星中较为常见，但并非普遍存在$^{[38]}$。此外，MIRI 还在行星质量伴星 VHS 1256b 中观测到了硅酸盐吸收$^{[39]}$。

作为大气对流过程一部分的“铁雨”只可能发生在褐矮星中，而不会出现在小质量恒星中。关于铁雨的光谱研究仍在进行中，但并非所有褐矮星在任何时候都会表现出这种大气异常。2013 年，人们在附近的 Luhman 16 系统中，对其 B 分量周围的非均匀含铁大气进行了成像$^{[40]}$。

对于晚期 T 型褐矮星，仅开展了少量的变光搜索研究。理论预测，在晚期 T 型褐矮星中会形成由铬、氯化钾以及多种硫化物构成的薄云层，这些硫化物包括硫化锰、硫化钠和硫化锌$^{[41]}$。对可变的 T7 型褐矮星 2M0050–3322 的解释认为，其最上层为氯化钾云，中间层为硫化钠云，下层为硫化锰云。上两层云的斑块状分布可以解释为何甲烷和水汽吸收带呈现出变化$^{[42]}$。

在最低温度的 Y 型褐矮星 WISE 0855−0714 中，硫化物云和水冰云的斑块状云层可能覆盖其约 50\% 的表面$^{[43]}$。
\subsubsection{低质量褐矮星与高质量行星的比较}
\begin{figure}[ht]
\centering
\includegraphics[width=6cm]{./figures/e55ea39cbe25dd2e.png}
\caption{} \label{fig_Brdw_7}
\end{figure}
与恒星类似，褐矮星也是独立形成的；但与恒星不同的是，它们缺乏足够的质量来“点燃”氢聚变。和所有恒星一样，褐矮星既可以单独存在，也可以与其他恒星近距离共存。有些褐矮星环绕恒星运行，并且像行星一样，其轨道也可能具有较大的偏心率。

\textbf{尺寸与燃料燃烧的不确定性}

褐矮星的半径大致都与木星相当。在其质量范围的高端（60–90 $M_{\mathrm{J}}$），褐矮星的体积主要受电子简并压控制，这一点与白矮星类似$^{[44]}$；而在质量范围的低端（约 10 $M_{\mathrm{J}}$），其体积主要受库仑压控制，这又与行星相似。综合来看，在整个可能的质量范围内，褐矮星的半径变化仅有约 10–15\%。此外，质量–半径关系在从约一个土星质量到氢燃烧起始点（$0.080\pm0.008,M_\odot$）之间几乎没有变化，这表明从这一角度看，褐矮星只是高质量的类木行星$^{[45]}$。这也使得在观测上区分褐矮星与行星变得十分困难。

另外，许多褐矮星根本不发生任何核聚变；即便是处在质量上限（超过 60 $M_{\mathrm{J}}$）的褐矮星，也会因冷却速度很快而在约一千万年后不再进行聚变反应。

\textbf{热辐射光谱}

X 射线和红外光谱是识别褐矮星的重要特征。一些褐矮星会发射 X 射线；而所有“温暖”的褐矮星都会在红光和红外波段持续发光，直到它们冷却到类似行星的温度（低于 1000 K）。

气态巨行星也具有某些与褐矮星相似的特性。与太阳类似，木星和土星主要由氢和氦组成。尽管土星的质量只有木星的约 30\%，其体积却几乎与木星相当。太阳系中的四颗巨行星中，有三颗（木星、土星和海王星）辐射出的热量远多于它们从太阳接收到的热量，最高可达约两倍$^{[46][47]}$。这四颗巨行星都拥有各自的“行星系统”，即由大量卫星构成的广泛卫星体系。

\textbf{现行 IAU 标准}

目前，国际天文学联合会（IAU）将质量高于 13 个木星质量（$13,M_{\mathrm{J}}$，即氘发生热核聚变的质量下限）的天体视为褐矮星，而将质量低于该数值（并且绕恒星或恒星遗迹运行）的天体视为行星。能够触发并维持氢燃烧所需的最小质量（约 $80,M_{\mathrm{J}}$）构成了这一分类的上限$^{[48]}$。

同时，学界也在讨论：与其依据基于核聚变反应的理论质量界限，不如依据形成过程来定义褐矮星是否更为合理$^{[3]}$。在这种解释下，褐矮星是恒星形成过程所产生的最低质量产物，而行星则是在恒星周围的吸积盘中形成的天体。已发现的一些最冷的自由漂浮天体（如 WISE 0855），以及已知质量最低的年轻天体（如 PSO J318.5−22），被认为其质量低于 $13,M_{\mathrm{J}}$，因此由于其究竟应被视为流浪行星还是褐矮星存在歧义，常被称为行星质量天体。此外，也已知存在绕褐矮星运行的行星质量天体，例如 2M1207b、2MASS J044144b 和 Oph 98 B。

将 $13,M_{\mathrm{J}}$ 作为分界线更多是一种经验法则，而非具有精确定义的物理常数。质量更大的天体会燃烧其大部分氘，而质量更小的天体只会燃烧极少量，$13,M_{\mathrm{J}}$ 大致位于两者之间$^{[49]}$。氘燃烧的程度还在一定程度上取决于天体的组成，尤其是氦和氘的含量，以及重元素所占比例；这些因素会影响大气的不透明度，从而影响辐射冷却速率$^{[50]}$。

截至 2011 年，《系外行星百科全书》（Extrasolar Planets Encyclopaedia）收录了质量高达 $25,M_{\mathrm{J}}$ 的天体，并指出：“在观测到的质量谱中，$13,M_{\mathrm{J}}$ 附近并不存在任何特殊特征，这一事实强化了放弃该质量界限的选择”$^{[51]}$。到 2016 年，基于质量–密度关系的研究，该上限被提高至 $60,M_{\mathrm{J}}$,$^{[52][53]}$。

Exoplanet Data Explorer 收录了质量最高达 $24,M_{\mathrm{J}}$ 的天体，并附有说明：“IAU 工作组提出的 $13,M_{\mathrm{J}}$ 区分标准在具有岩石核心的行星情形下缺乏物理动机，并且由于 $\sin i$ 不确定性，在观测上也存在问题”$^{[54]}$。NASA 系外行星档案（NASA Exoplanet Archive）则收录了质量（或最小质量）不超过 $30,M_{\mathrm{J}}$ 的天体$^{[55]}$。

\textbf{次褐矮星}
\begin{figure}[ht]
\centering
\includegraphics[width=6cm]{./figures/8c9a960c9abbf3d5.png}
\caption{太阳、一个年轻的次褐矮星与木星的体积大小对比。随着次褐矮星逐渐老化，它会不断冷却并收缩。} \label{fig_Brdw_8}
\end{figure}
质量低于 13 MJ 的天体被称为次褐矮星（sub-brown dwarfs）或行星质量褐矮星。它们的形成方式与恒星和褐矮星相同（即通过气体云的引力坍缩），但其质量低于氘发生热核聚变的下限 $^{[56]}$。有些研究者将它们称为自由漂浮行星$^{[57]}$，而另一些则称其为行星质量褐矮星 $^{[58]}$。
\subsubsection{其他物理性质在质量估计中的作用}
尽管光谱特征可以帮助区分低质量恒星与褐矮星，但在许多情况下仍需要通过质量估计才能得出结论。质量估计背后的理论认为，质量相近的褐矮星具有相似的形成过程，并且在形成初期温度较高。有些天体的光谱型与低质量恒星相似，例如 2M1101AB。随着时间推移而冷却，这些褐矮星会根据其质量保留不同的光度范围 $^{[59]}$。

如果缺乏年龄和光度信息，质量估计将变得十分困难。例如，一个 L 型褐矮星可能是质量较高、年龄较大的褐矮星（甚至可能是低质量恒星），也可能是质量极低、但非常年轻的褐矮星。对于 Y 型褐矮星而言，这个问题相对较小，因为即使在较大的年龄估计下，它们仍然是接近次褐矮星质量下限的低质量天体 $^{[60]}$。而对于 L 型和 T 型褐矮星，获得准确的年龄估计仍然十分重要。相比之下，光度并不是最棘手的问题，因为可以通过光谱能量分布来估计 $^{[61]}$。

年龄估计通常有两种方法：其一，褐矮星本身非常年轻，并且仍然保留与年轻阶段相关的光谱特征；其二，褐矮星与某颗恒星或恒星群（如疏散星团或恒星协会）具有共同运动特性，在这些系统中年龄更容易确定。

利用第二种方法研究的一个非常年轻的褐矮星实例是 2M1207 及其伴星 2M1207b。根据其空间位置、自行和光谱特征，该系统被确定属于约 800 万年历史的 TW Hydrae 恒星协会，其次级天体的质量被估计为$8 \pm 2M_J$，低于氘燃烧下限 $^{[62]}$。

而通过共同运动方法得到极高年龄估计的一个例子是褐矮星–白矮星双星系统 COCONUTS-1，其中白矮星的年龄估计为$7.3(^{+2.8}_{-1.6})$十亿年。在这一案例中，质量并非直接由年龄推导，而是通过共同运动获得了精确的距离估计（利用 Gaia 的视差测量）。在此基础上，研究者估计了该褐矮星的半径，并进一步推算其质量为$15.4(^{+0.9}_{-0.8}M_J)$.$^{[63]}$。
\subsection{观测}
\subsubsection{褐矮星的分类}
\textbf{光谱型 M}
\begin{figure}[ht]
\centering
\includegraphics[width=6cm]{./figures/0a732b5953b02bea.png}
\caption{晚期$M$型矮星的艺术想象图} \label{fig_Brdw_9}
\end{figure}
这些是光谱型为 M5.5 或更晚的棕矮星，也被称为晚期$M$型矮星。所有光谱型为 M 的棕矮星都是年轻天体，例如 Teide 1 ——这是发现的第一颗$M$型棕矮星——以及 LP 944-20，它是距离最近的$M$型棕矮星。

\textbf{光谱级 L}
\begin{figure}[ht]
\centering
\includegraphics[width=6cm]{./figures/a812750df0497cb6.png}
\caption{$L$型矮星的艺术概念图} \label{fig_Brdw_10}
\end{figure}
M 型光谱级（经典恒星序列中最冷的一类）的决定性特征，是其光学光谱由二氧化钛（TiO）和二氧化钒（VO）分子的吸收带所主导。然而，白矮星 GD 165 的低温伴星 GD 165B 却完全不具备 M 型矮星典型的 TiO 特征。随后，人们又发现了许多与 GD 165B 类似的天体，这最终促成了一个新的光谱类型——L 型矮星的确立。

L 型矮星在红光光学波段的定义特征，不再是金属氧化物（TiO、VO）的吸收带，而是金属氢化物（FeH、CrH、MgH、CaH）的发射带，以及碱金属原子（Na、K、Rb、Cs）显著的原子谱线。截至 2013 年，已确认的 L 型矮星超过 900 个$^{[26]}$，其中大多数是通过大视场巡天发现的，包括二维微米全天巡天（2MASS）、南天深近红外巡天（DENIS）以及斯隆数字化巡天（SDSS）。

这一光谱类型还包含了最冷的主序恒星（质量大于约 80 MJ），它们的光谱型介于 L2 到 L6 之间$^{[64]}$。

\textbf{光谱级 T}
\begin{figure}[ht]
\centering
\includegraphics[width=6cm]{./figures/b57b3066df5ac4fb.png}
\caption{$T$型矮星的艺术构想图} \label{fig_Brdw_11}
\end{figure}
由于 GD 165B 是 L 型矮星的原型，Gliese 229B 则成为第二个新光谱类型—— T 型矮星的原型。T 型矮星呈粉紫色（洋红色）外观。与 L 型矮星在近红外（NIR）光谱中表现出强烈的水$H(_2)O$和一氧化碳（CO）吸收带不同，Gliese 229B 的近红外光谱主要由甲烷$CH(_4)$吸收带主导；这种特征在太阳系中仅存在于气态巨行星及土卫六（Titan）的大气中。$CH(_4)$、$H(_2)O$ 以及分子氢（$H(_2)$）的碰撞诱导吸收（CIA）共同赋予 Gliese 229B 蓝色的近红外颜色。其红光学波段的光谱陡峭倾斜，并且缺乏$L$型矮星所特有的 FeH 和 CrH 吸收带，而是受到碱金属钠（Na）和钾（K）异常宽广吸收特征的显著影响。

正是由于这些差异，J. Davy Kirkpatrick 提出了$T$光谱型，用于描述在$H$波段和$K$波段中表现出甲烷$CH(_4)$吸收的天体。截至 2013 年，已知的$T$型矮星共有 355 个$^{[26]}$。近年来，Adam Burgasser 和 Tom Geballe 建立了$T$型矮星的近红外光谱分类体系。理论研究表明，$L$型矮星包含极低质量恒星与次恒星天体（褐矮星）的混合体，而$T$型矮星则完全由褐矮星组成。由于$T$型矮星在光谱绿色区域中对钠和钾的强烈吸收，其在人眼视觉下的实际颜色预计并非棕色，而是洋红色$^{[65][66]}$。

早期观测限制了$T$型矮星可被探测的最大距离。然而，如 WISE 0316+4307 这样的$T$型褐矮星，已在距离太阳超过 100 光年的位置被发现。利用 詹姆斯·韦布空间望远镜（JWST）的观测，人们甚至探测到了距离太阳高达 4500 秒差距的$T$型矮星，例如 UNCOVER-BD-1。

\textbf{光谱类型 Y}
\begin{figure}[ht]
\centering
\includegraphics[width=6cm]{./figures/02ecc07011945347.png}
\caption{$Y$型矮星的艺术想象图} \label{fig_Brdw_12}
\end{figure}
2009 年时，已知温度最低的棕矮星其有效温度估计在 500–600 K（227–327 °C；440–620 °F）之间，并被归入光谱型 T9。其代表性例子包括棕矮星 CFBDS J005910.90–011401.3、ULAS J133553.45+113005.2 以及 ULAS J003402.77−005206.7$^{[67]}$。这些天体的光谱在约 1.55 微米处显示出吸收峰$^{[67]}$。Delorme 等人提出，该特征源于氨的吸收，并认为这可作为$T$型向$Y$型过渡的标志，从而将这些天体归类为 Y0 型$^{[67][68]}$。然而，这一特征很难与水和甲烷的吸收区分开来$^{[67]}$，也有其他作者指出，将其直接定为 Y0 型仍为时过早$^{[69]}$。
\begin{figure}[ht]
\centering
\includegraphics[width=6cm]{./figures/0bd389856fb9037f.png}
\caption{WISE 0458+6434 是由 WISE 发现的第一颗超冷棕矮星（绿色圆点）。其中呈现的绿色和蓝色来自红外波段，这些红外波长被映射为可见颜色。} \label{fig_Brdw_13}
\end{figure}
\textbf{提出的光谱类型 H}

2025 年，天文学家 Kevin Luhman 和 Catarina Alves de Oliveira 提出了一种新的光谱类型$H$（$H$源自 hydrocarbon，即“烃类”）。他们利用詹姆斯·韦布空间望远镜的数据，在恒星形成区 IC 348 中识别出许多棕矮星，这些天体的质量非常低（其中许多低于氘燃烧极限），并且在 3.4 μm 处呈现出一条吸收线，该吸收线对应于一种尚未被明确识别的脂肪族烃。这种吸收线将作为$H$型光谱的定义特征。$^{[71]}$

\textbf{垂直混合的作用}
\begin{figure}[ht]
\centering
\includegraphics[width=6cm]{./figures/2efdd4bc98a101d2.png}
\caption{连接一氧化碳与甲烷的主要化学反应路径。短寿命自由基以圆点标注。改编自 Zahnle 与 Marley。$^{[72]}$} \label{fig_Brdw_14}
\end{figure}
在以氢为主的大气中，棕矮星内部存在着一氧化碳（CO）与甲烷（CH₄）之间的化学平衡。在正向反应中，一氧化碳与氢分子反应，生成甲烷和羟基自由基（OH）。随后，羟基自由基可能再与氢反应，形成水分子。在反向反应中，甲烷与羟基反应，重新生成一氧化碳和氢。

在较高温度和较低压力条件下（对应$L$型矮星），该化学反应更倾向于一氧化碳；而在较低温度和较高压力条件下（对应$T$型矮星），反应则更倾向于甲烷，在$T/Y$边界附近由甲烷占主导地位。然而，大气中的垂直混合会改变这种平衡：甲烷可能下沉到更深层的大气，而一氧化碳则会从更深、更高温的层次上升到上部大气。由于 C–O 键的断裂存在能垒，一氧化碳重新转化为甲烷的反应速度较慢，这使得棕矮星的可观测大气处于一种化学非平衡状态。

因此，$L/T$转变主要被定义为：从$L$型矮星中以一氧化碳为主的大气，过渡到$T$型矮星中以甲烷为主的大气。垂直混合的强弱会将这一转变推向更低或更高的温度。这一点对于表面重力较小、且具有扩展大气层的天体尤为重要，例如巨型系外行星。在这些行星中，$L/T$ 转变会发生在更低的温度下；而对于棕矮星而言，这一转变通常发生在约 1200 K 左右。相比之下，系外行星 HR 8799c 的温度约为 1100 K，却并未显示出甲烷特征。$^{[72]}$

$T$型与$Y$型矮星之间的转变通常被定义在约 500 K 左右，这主要是由于这些寒冷且极其昏暗的天体缺乏足够的光谱观测数据。未来借助詹姆斯·韦布空间望远镜（JWST）以及超大望远镜（ELTs）的观测，有望显著扩展具有光谱观测的$Y$型矮星样本。$Y$型矮星的大气通常由强烈的甲烷和水汽吸收特征所主导，并可能还存在氨以及水冰的吸收特征。$^{[73]}$

此外，垂直混合、云层结构、金属丰度、光化学过程、闪电、撞击激波以及金属催化作用等因素，都可能影响$L/T$转变和$T/Y$转变所发生的具体温度。$^{[72]}$

\textbf{次级特征}

年轻的棕矮星具有较低的表面重力，这是因为与具有相似光谱型的场星相比，它们半径更大而质量更低。这类天体通常用后缀字母来标注其表面重力特征：$\beta$（beta）表示中等表面重力，γ（gamma）表示低表面重力。低表面重力的判据包括：CaH、K I 和 Na I 吸收线较弱，以及VO 吸收线较强。$^{[76]}$α（alpha）表示正常表面重力，通常会被省略不写。有时，极低的表面重力会用δ（delta）来表示。$^{[78]}$后缀 “pec”表示 peculiar（异常），该后缀仍用于描述其他不寻常的特征，通常是对多种性质的综合概括，可指示低表面重力、次矮星（subdwarf）或未分辨的双星系统等情况。$^{[79]}$前缀 “sd”表示次矮星，仅用于冷次矮星，该前缀意味着低金属丰度，并且其运动学性质更接近银河晕星而非银盘星。$^{[75]}$ 次矮星在颜色上通常比银盘天体更偏蓝。$^{[80]}$红色后缀用于描述颜色偏红的天体，这通常与较老的年龄有关，而不被解释为低表面重力，而是与较高的尘埃含量有关。$^{[77]}$ $^{[78]}$蓝色后缀则描述近红外颜色偏蓝、且不能用低金属丰度来解释的天体。其中一部分被解释为$L+T$型双星系统，而另一些并非双星，例如 2MASS J11263991−5003550，其性质被认为与稀薄和/或大颗粒云层有关。$^{[78]}$
\subsubsection{棕矮星的光谱与大气性质}
\begin{figure}[ht]
\centering
\includegraphics[width=6cm]{./figures/ebca80dcfaaa1db9.png}
\caption{棕矮星内部结构的艺术示意图。由于层状位移，特定深度处的云层出现错位。} \label{fig_Brdw_15}
\end{figure}
$L$型和$T$型矮星所辐射的绝大部分能量通量集中在 1–2.5 微米的近红外波段。沿着晚期$M$型、$L$型到$T$型矮星序列，温度不断降低并持续下降，这使得近红外光谱极为丰富，包含多种多样的谱线特征：既有相对狭窄的中性原子谱线，也有宽广的分子吸收带，而这些特征对温度、表面重力以及金属丰度的依赖关系各不相同。此外，这种低温环境有利于气体发生凝结并形成固态颗粒。
\begin{figure}[ht]
\centering
\includegraphics[width=6cm]{./figures/6c041779a34cc547.png}
\caption{测量到的风（斯皮策空间望远镜；艺术概念图；2020 年 4 月 9 日）$^{[81]}$} \label{fig_Brdw_16}
\end{figure}
已知褐矮星的典型大气温度范围约为 2200 K 至 750 K$^{[65]}$。与依靠稳定内部核聚变自我加热的恒星不同，褐矮星会随时间迅速冷却；质量较大的褐矮星冷却得比质量较小的更慢。有一些证据表明，在光谱型$L$与$T$的过渡区（约 1000 K）附近，褐矮星的冷却过程会减缓$^{[82]}$。

对已知褐矮星候选体的观测揭示了一种红外辐射周期性变亮与变暗的模式，这表明相对较冷且不透明的云层结构遮蔽着在强烈风场搅动下的炽热内部。这类天体上的“天气”被认为极其猛烈，其强度可与、甚至远超木星著名的风暴。

2013 年 1 月 8 日，天文学家利用 NASA 的哈勃空间望远镜和斯皮策空间望远镜，对一颗名为 2MASS J22282889–4310262 的褐矮星的多风暴大气进行了探测，绘制出了迄今为止最为详细的褐矮星“天气图”。该图显示了由风驱动的、行星尺度的云层。上述研究为更深入理解褐矮星以及太阳系外行星的大气性质迈出了重要一步$^{[83]}$。

2020 年 4 月，科学家报告称在附近的一颗褐矮星 2MASS J10475385+2124234 上测得风速为 $+650 \pm 310$ m/s（最高可达约 1450 英里/小时）。为获得这一结果，研究人员将由亮度变化推断出的气象特征旋转运动，与褐矮星内部产生的电磁旋转进行对比。该结果验证了此前关于褐矮星具有高速风场的预测。科学家们期望，这种对比方法可用于研究其他褐矮星及系外行星的大气动力学$^{[84]}$。
\subsubsection{观测技术}
\begin{figure}[ht]
\centering
\includegraphics[width=6cm]{./figures/559b81ff99676a50.png}
\caption{褐矮星 Teide 1、Gliese 229B和WISE 1828+2650与红矮星 Gliese 229A、木星以及太阳的对比} \label{fig_Brdw_17}
\end{figure}
近来，日冕仪（coronagraph）已被用于探测绕明亮可见恒星运行的暗弱天体，其中包括 Gliese 229B。

配备电荷耦合器件（CCD）的高灵敏度望远镜被用于搜索遥远星团中的暗弱天体，例如 Teide 1。

大视场巡天已经识别出一些单个的暗弱天体，如 Kelu-1（距离地球约 30 光年）。

褐矮星常常是在系外行星巡天中被发现的。用于探测系外行星的方法同样适用于褐矮星，而且由于褐矮星本身更亮，它们实际上更容易被探测到。

由于具有强磁场，褐矮星可以成为强烈的射电辐射源。阿雷西博天文台和甚大阵列（VLA）的观测计划已经探测到十多个这样的天体，它们也被称为超冷矮星，因为它们与该类中其他天体共享相似的磁性质$^{[85]}$。对褐矮星射电辐射的探测使得其磁场强度可以被直接测量。
\subsubsection{里程碑}
\begin{itemize}
\item 1995 年：首个褐矮星得到确认。位于昴宿星团中的 Teide 1（一颗 M8 型天体）在西班牙加那利群岛天体物理研究所罗克·德洛斯·穆查乔斯天文台，借助 CCD 被识别出来。
首个甲烷褐矮星得到确认。Gliese 229B 被发现绕红矮星 Gliese 229A（距离约 20 光年）运行；该发现使用了自适应光学日冕仪，以增强加州南部帕洛玛山帕洛玛天文台 60 英寸（1.5 m）反射望远镜的成像效果；随后用 200 英寸（5.1 m）的 Hale 望远镜进行的近红外光谱观测显示其大气中富含甲烷。
\item 1998 年：发现首个发射 X 射线的褐矮星。位于蝘蜓座 I 暗云中的 Cha Helpha 1（M8 型天体）被确定为 X 射线源，其性质与具有对流外层的晚型恒星相似。
\item 1999 年 12 月 15 日：首次在褐矮星中探测到 X 射线耀斑。加州大学的一支团队利用 钱德拉 X 射线天文台监测 LP 944-20（质量约 60 MJ，距离约 16 光年），捕捉到一次持续约 2 小时的耀斑$^{[86]}$。
\item 2000 年 7 月 27 日：首次从褐矮星中探测到射电辐射（包括耀斑态与静态辐射）。甚大阵列（VLA）的一组学生团队探测到 LP 944–20 的射电辐射$^{[87]}$。
\item 2004 年 4 月 30 日：首次在褐矮星周围探测到系外行星候选体：使用 VLT 发现 2M1207b，同时这也是首颗被直接成像的系外行星$^{[88]}$。
\item 2013 年 3 月 20 日：发现距离最近的褐矮星系统：Luhman 16$^{[89]}$。
\item 2014 年 4 月 25 日：发现当时已知最冷的褐矮星。WISE 0855−0714距离地球约 7.2 光年（为距离太阳第七近的恒星系统），其温度介于 −48 ℃ 至 −13 ℃ 之间$^{[90]}$。
\end{itemize}
\subsubsection{褐矮星的 X 射线源}
\begin{figure}[ht]
\centering
\includegraphics[width=8cm]{./figures/0dcede50dee6ee4c.png}
\caption{LP 944-20 在耀斑发生前与耀斑期间的钱德拉（Chandra）X 射线图像} \label{fig_Brdw_18}
\end{figure}
自 1999 年以来在褐矮星中探测到的 X 射线耀斑表明，其内部磁场正在发生变化，类似于极低质量恒星中的磁场。尽管褐矮星不像恒星那样在核心中将氢聚变成氦，但氘聚变以及引力收缩所释放的能量仍能使其内部保持高温，并产生强磁场。褐矮星的内部处于一种快速翻滚、即强烈对流的状态。再加上大多数褐矮星具有的快速自转，这种对流为在其“表面”附近形成强烈而紊乱的磁场创造了条件。钱德拉（Chandra）望远镜在 LP 944-20 上观测到的耀斑，其产生的磁场正是起源于褐矮星“表面”之下的湍动磁化等离子体。

利用美国国家航空航天局（NASA）的钱德拉 X 射线天文台，科学家在一个多恒星系统中探测到了来自一颗低质量褐矮星的 X 射线信号$^{[91]}$。这是首次在 X 射线波段分辨出一颗如此接近其母恒星（类太阳恒星 TWA 5A）的褐矮星$^{[91]}$。东京中央大学的坪井洋子（Yohko Tsuboi）表示：“我们的钱德拉数据表明，这些 X 射线起源于褐矮星的日冕等离子体，其温度约为 300 万摄氏度$^{[91]}$。”坪井还指出：“这颗褐矮星在 X 射线波段的亮度与当今的太阳相当，而其质量却只有太阳的五十分之一$^{[91]}$。因此，这一观测提出了一种可能性：即便是质量巨大的行星，在其年轻阶段也可能自行辐射 X 射线！$^{[91]}$”
\subsubsection{作为射电源的褐矮星}
\begin{figure}[ht]
\centering
\includegraphics[width=6cm]{./figures/cb4bae58f97ccb85.png}
\caption{LSR J1835+3259 的静态射电辐射（等高线）被成功分辨出来，其形态与木星的辐射带相似。中央的暗斑是来自右旋圆偏振极光的射电辐射。} \label{fig_Brdw_19}
\end{figure}
最早被发现能够发射射电信号的棕矮星是 LP 944-20。之所以对它进行观测，是因为它同时也是一个 X 射线源，而这两类辐射都是日冕存在的标志。大约 5–10\% 的棕矮星似乎具有强磁场并能发射射电波；基于蒙特卡洛建模及其平均空间密度的估计，在距离太阳 25 pc以内，可能多达 40 颗具有磁性的棕矮星$^{[92]}$。棕矮星的射电辐射功率在不同温度之间变化不大，整体上近似保持恒定$^{[85]}$。棕矮星可能维持强度高达 6 kG 的磁场$^{[93]}$。天文学家根据射电辐射的性质估计，棕矮星的磁层高度可达约 $10^7m$,$^{[94]}$。目前尚不清楚棕矮星的射电辐射在本质上更接近行星还是恒星的射电辐射。一些棕矮星会发射规则的射电脉冲，这有时被解释为来自两极的定向射电辐射，但也可能源自活动区域的定向辐射。射电波偏振方向的规则、周期性反转，可能表明棕矮星的磁场会周期性地发生极性反转。这种反转可能源于类似太阳活动周期的棕矮星磁活动周期$^{[95]}$。

第一颗被发现能够发射射电波的$M$光谱型 棕矮星是 LP 944-20，于 2001 年探测到；第一颗$L$光谱型、能够发射射电波的棕矮星是 2MASS J0036159+182110，于 2008 年探测到；第一颗$T$光谱型、能够发射射电波的棕矮星是 2MASS J10475385+2124234$^{[96][97]}$。这一最后的发现尤为重要，因为它表明温度与系外行星相近的棕矮星也能够拥有强度大于 1.7 kG 的磁场。尽管在 2010 年阿雷西博天文台曾对 Y 光谱型棕矮星进行了高灵敏度的射电辐射搜索，但并未探测到射电信号$^{[98]}$。
\subsubsection{最新进展}
\begin{figure}[ht]
\centering
\includegraphics[width=8cm]{./figures/bb2e261b10e71e92.png}
\caption{一幅可视化图像，展示了在距离太阳 65 光年范围内已被发现的棕矮星（红点）的三维分布图$^{[99]}$} \label{fig_Brdw_20}
\end{figure}
对太阳邻域中棕矮星数量的估计表明，每出现一颗棕矮星，可能对应多达六颗恒星。$^{[100]}$而一项发表于 2017 年、基于年轻大质量恒星团 RCW 38 的更新研究则认为，银河系中可能包含约 250 亿到 1000 亿颗棕矮星。$^{[101]}$（可将这一数值与银河系恒星总数的估计作比较：约 1000 亿到 4000 亿颗。）

在 2017 年 8 月发表的一项研究中，NASA 的斯皮策空间望远镜监测了棕矮星红外亮度的变化，这些变化由厚度不均匀的云层覆盖所引起。观测结果揭示了在棕矮星大气中传播的大尺度波动（类似于海王星及太阳系其他巨行星的大气）。这些大气波会调制云层厚度，并以不同速度传播（可能源于差异自转）。$^{[102]}$

2020 年 8 月，天文学家通过“后院世界：第九行星”（Backyard Worlds: Planet 9）项目，在太阳附近发现了 95 颗棕矮星。$^{[103]}$

2024 年，詹姆斯·韦布空间望远镜（JWST）对两颗棕矮星给出了迄今为止最为详细的“天气报告”，揭示了其风暴频发的大气环境。这两颗棕矮星属于一个在 2013 年发现的双星系统 Luhman 16，距离地球仅 6.5 光年，是距离太阳最近的棕矮星。研究人员发现，它们具有高度湍动的云层，云粒很可能由硅酸盐颗粒构成，温度范围约为 875 °C（1607 °F）至 1026 °C（1879 °F），表明在这些棕矮星上存在被强风吹动的“高温沙尘”。此外，还探测到了一氧化碳、甲烷和水蒸气的吸收特征。$^{[104]}$
\subsection{双棕矮星系统}
\subsubsection{棕矮星—棕矮星双星系统}
\begin{figure}[ht]
\centering
\includegraphics[width=8cm]{./figures/8b70f1da37ff9cfb.png}
\caption{使用哈勃空间望远镜在多个历元拍摄的棕矮星双星图像。其中，双星系统 Luhman 16 AB（左）比此处展示的其他示例更接近太阳系。} \label{fig_Brdw_21}
\end{figure}
M、L 和 T 型的棕矮星双星在主星质量越低时越不常见$^{[105]}$。L 型棕矮星的双星比例约为 (24^{+6}_{-2}\%)，而晚期 T 型—早期 Y 型棕矮星（T5–Y0）的双星比例约为 (8\%\pm6\%)$^{[106]}$。

与高质量体系相比，棕矮星双星在低质量体系中具有更高的伴星—主星质量比$q=\frac{M_B}{M_A}$.例如，以 M 型恒星为主星的双星，其$q$的分布范围较宽，但偏好$q\ge 0.4$。相比之下，棕矮星双星则明显偏好$q\ge 0.7$。双星分离度随质量降低而减小：$M$型恒星的分离度峰值在 3–30 个天文单位（au）；M–L 型棕矮星的投影分离度峰值在 5–8 au；而 T5–Y0 型天体的投影分离度服从对数正态分布，其峰值约为 2.9 au$^{[106]}$。

一个典型例子是距离最近的棕矮星双星 Luhman 16 AB，其主星为 L7.5 型棕矮星，分离度为 3.5 au，质量比 $q=0.85$。这一分离度位于 M–L 型棕矮星预期分离范围的低端，但其质量比是典型的。

目前尚不清楚这一趋势是否会延续到$Y$型棕矮星，因为相关样本数量极少。Y+Y 型棕矮星双星预计具有较高的质量比$q$和较小的分离度，甚至可小于 1 au$^{[107]}$。2023 年，Y+Y 型棕矮星 WISE J0336−0143 通过 JWST 观测被确认为双星，其质量比为 (q=0.62\pm0.05)，分离度为 0.97 au。研究人员指出，目前低质量棕矮星双星的样本规模仍不足以判断 WISE J0336−0143 是低质量双星的典型代表，还是一个特例系统$^{[108]}$。

通过观测包含棕矮星的双星系统轨道，可以测量棕矮星的质量。在 2MASSW J0746425+2000321 系统中，次星的质量约为太阳质量的 6\%。这种测量被称为动力学质量$^{[109][110]}$。距离太阳系最近的棕矮星系统是双星 Luhman 16。研究人员曾尝试用类似的方法在该系统中搜索行星，但并未发现任何行星$^{[111]}$。
\subsubsection{异常的棕矮星双星}
\begin{figure}[ht]
\centering
\includegraphics[width=6cm]{./figures/8c1d53b83e52c057.png}
\caption{宽分离棕矮星双星 SDSS J1416+1348} \label{fig_Brdw_22}
\end{figure}
宽双星系统 2M1101AB 是首个分离距离大于 20 AU 的棕矮星双星。该系统的发现为理解棕矮星的形成机制提供了决定性的线索。此前人们认为，宽分离的棕矮星双星要么无法形成，要么会在 1–10 Myr 的年龄尺度内被破坏。该系统的存在也与抛射假说不一致$^{[112]}$。抛射假说认为，棕矮星形成于多星系统中，但在获得足够质量以点燃氢燃烧之前就被抛射出去$^{[113]}$。

近些年又发现了宽双星 W2150AB。它与 2M1101AB 具有相似的质量比和束缚能，但年龄更大，且位于银河系的不同区域。2M1101AB 位于恒星密集区域，而 W2150AB 则处在较为稀疏的恒星场中，因此它必须在其诞生的星团中幸存过各种动力学相互作用。该系统也属于少数可以由地基望远镜轻易分辨的 L+T 型双星之一；另外两个是 SDSS J1416+13AB 和 Luhman 16$^{[114]}$。

还有一些其他有趣的双星系统，例如食双星棕矮星系统 2MASS J05352184–0546085$^{[115]}$。对该系统的光度学研究显示，其中质量较小的棕矮星反而比质量较大的伴星更热$^{[116]}$。
\subsubsection{围绕恒星运行的棕矮星}
在恒星周围的近距离轨道（小于 5 AU）上同时存在棕矮星和大质量行星的情况十分罕见，这种现象通常被称为“棕矮星荒漠”。质量与太阳相当的恒星中，只有不到 1\% 在 3–5 AU 范围内拥有一颗棕矮星$^{[117]}$。

恒星—棕矮星双星的一个例子是最早发现的$T$型棕矮星 Gliese 229 B，它绕主序红矮星 Gliese 229 A 运行。也已发现围绕亚巨星运行的棕矮星，例如 TOI-1994b，其轨道周期为 4.03 天$^{[118]}$。

关于某些低质量棕矮星是否应被视为行星，学界仍存在分歧。NASA 系外行星档案库将最小质量不超过 30 个木星质量、且满足其他条件（如绕恒星运行）的天体也列为行星$^{[119]}$；而国际天文学联合会（IAU）的系外行星工作组（WGESP）则只将质量低于 13 个木星质量的天体视为行星$^{[120]}$。
\subsubsection{白矮星—棕矮星双星}
\begin{figure}[ht]
\centering
\includegraphics[width=6cm]{./figures/920b1ba8f6661b98.png}
\caption{LSPM J0241+2553AB：一个宽分离的白矮星（A）—棕矮星（B）双星系统。} \label{fig_Brdw_23}
\end{figure}
围绕白矮星运行的棕矮星相当罕见。GD 165 B 是其中一个典型系统，它也是 L 型棕矮星的原型 $^{[121]}$。这类系统在确定体系年龄以及棕矮星质量方面非常有价值。其他白矮星—棕矮星双星系统还包括 COCONUTS-1 AB（约 70 亿年）$^{[63]}$、LSPM J0055+5948 AB（约 100 亿年）$^{[122]}$、SDSS J22255+0016 AB（约 20 亿年）$^{[123]}$，以及 WD 0806−661 AB（约 15–27 亿年）$^{[124]}$。

那些具有近距离、潮汐锁定的棕矮星绕白矮星运行的系统，被归类为“公共包层后双星”（post-common-envelope binaries，PCEBs）。目前已确认的、包含白矮星与棕矮星伴星的 PCEB 仅有 8 个，其中包括 WD 0137-349 AB。在这些紧密的白矮星—棕矮星双星的演化历史中，棕矮星曾在母星处于红巨星阶段时被其包吞。质量低于 20 个木星质量的棕矮星会在这一包吞过程中被蒸发掉 $^{[125][126]}$。围绕白矮星近距离运行的棕矮星稀缺现象，可与主序星周围棕矮星的类似观测结果进行比较，这种现象通常被称为“棕矮星荒漠” $^{[127][128]}$。PCEB 体系可能进一步演化为灾变变星（CV*），其中棕矮星充当物质供体 $^{[129]}$。数值模拟表明，高度演化的 CV 大多与亚恒星质量的供体有关（比例可达 80\%）$^{[130]}$。一种称为 WZ Sge 型矮新星的 CV，其供体质量往往接近低质量恒星与棕矮星的分界线 $^{[131]}$。双星 BW Sculptoris 就是这样一个以棕矮星为供体的矮新星系统。该棕矮星很可能是在供体恒星失去足够质量后形成的。质量损失伴随着轨道周期的缩短，直到周期降至约 70–80 分钟的最小值，随后周期再次增长，这一演化阶段因此被称为“周期反弹者”（period bouncer）$^{[130]}$。此外，也可能存在棕矮星与白矮星并合的情形，新星 CK Vulpeculae 就可能是这样一次白矮星—棕矮星并合的产物 $^{[132][133]}$。
\subsection{形成与演化}
\begin{figure}[ht]
\centering
\includegraphics[width=6cm]{./figures/1b0575e65d9fc67e.png}
\caption{由位于$\delta$猎户座星团外围的棕矮星 Mayrit 1701117发射出的 HH 1165 喷流} \label{fig_Brdw_24}
\end{figure}
棕矮星形成的最早阶段被称为原棕矮星（proto-brown dwarf）或前棕矮星（pre-brown dwarf）。原棕矮星是原恒星（0 / I 类天体）的低质量对应物。此外，极低光度天体（Very Low Luminosity Objects，VeLLOs，满足 (L_{\mathrm{int}} \le 0.1\text{–}0.2,L_{\odot})）往往也是原棕矮星。它们通常分布在附近的恒星形成云中。截至 2024 年，已知约有 67 个较有希望的原棕矮星候选体和 26 个前棕矮星 $^{[134]}$。截至 2017 年，仅有一个已知与大型赫比格–哈罗（Herbig–Haro）天体相关联的原棕矮星，即 Mayrit 1701117；该棕矮星被一个伪盘和一个开普勒盘所环绕 $^{[135]}$。Mayrit 1701117 会喷射出长度约 0.7 光年的 HH 1165 喷流，该喷流主要在电离硫谱线中可见 $^{[136][137]}$。

棕矮星的形成方式与恒星相似，其周围同样存在原行星盘 $^{[138]}$，例如 Cha 110913−773444。研究发现，棕矮星周围的盘在许多方面与恒星周围的盘相似，因此人们预计在棕矮星周围也会存在通过吸积形成的行星 $^{[138]}$。鉴于棕矮星盘的质量较小，这些行星大多应为类地行星，而非气态巨行星 $^{[139]}$。如果一颗巨行星沿我们的视线方向掠过一颗棕矮星轨道，由于两者直径大致相当，通过凌日方法将产生很强的探测信号 $^{[140]}$。不过，棕矮星周围行星的吸积区距离棕矮星本身非常近，因此潮汐力将产生显著影响 $^{[139]}$。
\begin{figure}[ht]
\centering
\includegraphics[width=6cm]{./figures/f9bb69ab419b2bbf.png}
\caption{棕矮星 W1200-7845 的艺术家想象图} \label{fig_Brdw_25}
\end{figure}
2020 年，Disk Detective 项目发现了距离最近、仍保有原始盘（II 类盘）的棕矮星—— WISEA J120037.79−784508.3（W1200-7845）。这一发现源于分类志愿者注意到其红外过量辐射。随后，该天体经过科研团队的审核与分析，结果表明 W1200-7845 以 99.8\% 的概率属于 ε 船底座（ε Chamaeleontis，ε Cha）年轻运动群协会。基于 Gaia DR2 数据得到的视差测量，其距离地球约 102 秒差距（约 333 光年），位于本地太阳邻域之内 $^{[141][142]}$。
\begin{figure}[ht]
\centering
\includegraphics[width=6cm]{./figures/e45c6a97bfb0729e.png}
\caption{使用哈勃望远镜（Hubble）和詹姆斯·韦布空间望远镜（JWST）观测到的猎户座星云中的棕矮星原行星盘（proplyds）。} \label{fig_Brdw_26}
\end{figure}
2021 年的一项研究考察了年龄为数百万年、距离约 140–200 秒差距的恒星协会中，棕矮星周围的环星盘。研究人员发现，这些盘的质量不足以在未来形成行星。不过，在这些盘中也发现了一些迹象，可能表明行星形成在更早阶段就已开始，且行星可能已经存在于盘中。支持盘演化的证据包括：盘质量随时间减小、尘埃颗粒长大以及尘埃沉降$^{[143]}$。此外，还发现有两处棕矮星盘以吸收形式出现，并且在猎户座星云中至少有 4 个盘正在受到外部紫外辐射的光致蒸发；这类天体也被称为 proplyd（原行星盘）。其中，Proplyd 181−247（一颗棕矮星或低质量恒星）被一个半径约 30 天文单位的盘所环绕，盘质量为 6.2±1.0 个木星质量$^{[144]}$。
通常，棕矮星周围的盘半径小于 40 天文单位；但在更遥远的金牛座分子云中，有 3 个盘的半径大于 70 天文单位，并已被 ALMA 分辨成像。这些更大的盘具备形成质量大于 1 个地球质量的岩质行星的能力 $^{[145]}$。此外，在年龄超过数百万年的协会中也发现了仍保有盘的棕矮星$^{[146]}$，这可能表明棕矮星周围的盘需要更长时间才能消散。其中，特别古老的盘（>20 Myr）有时被称为“彼得·潘盘”（Peter Pan disks）。目前，2MASS J02265658−5327032是唯一已知拥有彼得·潘盘的棕矮星 $^{[147]}$。

位于变色龙座、距离地球约 500 光年的棕矮星 Cha 110913−773444，可能正处于形成一个微型行星系统的过程中。宾夕法尼亚州立大学的天文学家探测到一个气体与尘埃盘，其性质与被认为形成太阳系的原始盘相似。Cha 110913−773444是迄今发现质量最小的棕矮星（约 8 个木星质量）；如果它真的形成了行星系统，那么它将成为已知拥有行星系统的最小质量天体 $^{[148]}$。
\subsection{棕矮星周围的行星}
\begin{figure}[ht]
\centering
\includegraphics[width=6cm]{./figures/ea1832479cd56218.png}
\caption{围绕棕矮星的尘埃与气体盘的艺术想象图$^{[149]}$} \label{fig_Brdw_27}
\end{figure}
根据国际天文学联合会（IAU）在 2018 年 8 月给出的工作定义，系外行星可以绕棕矮星运行。其要求是：行星质量小于 13 个木星质量（MJ），且质量比满足$M/M_{\mathrm{central}} < \frac{2}{25+\sqrt{621}}$,约等于 1/25。这意味着：一颗质量为 80 MJ 的棕矮星，其周围质量不超过 3.2 MJ 的天体可被视为行星；而一颗质量为 13 MJ 的棕矮星，其周围质量不超过 0.52 MJ 的天体同样可被视为行星 $^{[150]}$。

围绕棕矮星、位于较大轨道距离处运行的“超木星”级行星质量天体，如 2M1207b、2MASS J044144 和 Oph 98 B，可能是通过分子云塌缩而非吸积形成的，因此更可能属于次棕矮星而非真正意义上的行星；这一判断依据其相对较大的质量和轨道半径。利用径向速度技术在较小轨道距离上发现首个绕棕矮星运行的低质量伴星（ChaHα8）为在数个天文单位甚至更小轨道尺度上探测棕矮星行星铺平了道路 $^{[151][152]}$。然而，由于 ChaHα8 系统中伴星与主星的质量比约为 0.3，该系统更类似于双星系统。随后在 2008 年，人们首次发现了一颗在相对较小轨道上绕棕矮星运行的行星质量伴星——MOA-2007-BLG-192Lb$^{[153]}$。

绕棕矮星运行的行星很可能是贫水的碳行星 $^{[154]}$。

一项基于 Spitzer 观测的 2017 年研究估计，为了以 95\% 的置信度通过凌日法至少探测到一颗小于地球尺寸的行星，需要监测约 175 颗棕矮星 $^{[155]}$。詹姆斯·韦布空间望远镜（JWST）有望探测到更小的行星。太阳系中行星与卫星的轨道往往与其所绕转的中心天体自转轴方向一致。若假设绕棕矮星或行星质量天体运行的行星轨道与其自转轴对齐，则类似木卫一（Io）的天体，其几何凌日概率可由公式 (\cos(79.5^\circ)/\cos(\text{倾角})) 计算 $^{[156]}$。多颗棕矮星和行星质量天体的轨道倾角已被估计，例如 SIMP 0136 的倾角约为 (80^\circ \pm 12^\circ) $^{[157]}$。若取 SIMP 0136 的下限倾角 (i \ge 68^\circ)，则其近轨行星的凌日概率可达 ≥48.6\%。不过，目前尚不清楚近轨行星在棕矮星周围的普遍程度；鉴于原行星盘尺度随质量减小而缩小，近轨行星可能在低质量天体周围更为常见 $^{[143]}$。

2025 年，人们提出了围绕 2M1510 的一颗极轨道双星行星的强有力证据，该发现由甚大望远镜（VLT）完成 $^{[158][159]}$。
\subsubsection{宜居性}
针对假想的绕棕矮星运行的行星的宜居性已有研究。计算机模型表明，这类行星具备宜居条件的限制极为严格：宜居带非常狭窄、距离中心天体极近（T 型棕矮星约为 0.005 天文单位），且由于棕矮星随时间冷却（其核聚变阶段最多仅持续约一千万年），宜居带会不断向内收缩。行星轨道必须具有极低的偏心率（约 (10^{-6}) 量级），以避免强烈的潮汐作用引发失控温室效应，从而使行星变得不宜居。此外，这类行星也不会拥有卫星 $^{[160]}$。
\subsection{极端棕矮星}
早在 1984 年，一些天文学家曾提出假设，认为太阳可能被一颗尚未探测到的棕矮星（有时称为 Nemesis）所环绕，该天体可能像掠过的恒星一样扰动奥尔特云。然而，这一假说如今已基本被放弃 $^{[161]}$。
\subsubsection{首次发现一览表}
\begin{table}[ht]
\centering
\caption{请输入表格标题}\label{tab_Brdw1}
\begin{tabular}{|c|c|c|c|c|c|}
\hline
\textbf{记录}& \textbf{名称} & \textbf{光谱型} & \textbf{赤经/赤纬} & \textbf{星座} & \textbf{备注} \\
\hline
首次探测到的 & Gliese 569 Bab（M3 型场星的伴星）& M8.5 与 M9& 14h54m29.2s，+16°06′04″ & 牧夫座（Boötes）& 1985 年成像，1988 年发表，1999 年分辨为双星伴星，2004 年测得质量\\
\hline 
首个场星型（孤立）且首个经确认的直接成像对象 & Teide 1 & M8 & 3h47m18.0s，+24°22′31″& 金牛座（Taurus）& 1994 年探测到，1995 年确认\\
\hline
首个具有晚 M 型光谱的对象 & Teide 1 & M8 & 3h47m18.0s，+24°22′31″ & 金牛座（Taurus） & 1994 年探测到\\
\hline
首个使用日冕仪成像且首个作为恒星伴星的对象 & Gliese 229 B & T6.5 & 06h10m34.62s，−21°51′52.1″ & 天兔座（Lepus） & 1994 年探测到，1995 年确认\\
\hline
首个具有 T 型光谱的对象 & Gliese 229 B & T6.5 & 06h10m34.62s，−21°51′52.1″ & 天兔座（Lepus） & 1994 年探测到\\
\hline
首个具有 L 型光谱的对象 & GD 165B & L4 & 14h24m39.144s，+09°17′13.98″ & 牧夫座（Boötes） & 1988 年发表，1999 年确认\\
\hline
首个三重棕矮星系统 & DENIS-P J020529.0−115925 A/B/C & L5、L8 与 T0 & 02h05m29.40s，−11°59′29.7″ & 鲸鱼座（Cetus） & Delfosse 等，1997 年 $^{[162]}$\\
\hline
首次发射 X 射线的棕矮星 & ChaHα1 & M8 & &变色龙座（Chamaeleon）&1998 年\\
\hline
首次观测到 X 射线耀斑 & LP 944–20 & M9V & 03h39m35.22s，−35°25′44.1″ & 天炉座（Fornax）& 1999 年\\
\hline
首个光谱双星棕矮星 & PPL 15 A、B$^{[163]}$ & M6.5 & 03h48m04.659s，+23°39′30.32″ & 金牛座（Taurus）& Basri 与 Martín，1999 年\\
\hline
首次探测到射电辐射（耀斑态与静态）& LP 944–20 & M9V & 03h39m35.22s，−35°25′44.1″ & 天炉座（Fornax）& 2000 年 $^{[87]}$\\
\hline
首个具有环星盘的棕矮星 & ChaHα1 & M7.5 & 11h07m17.0s，−77°35′54″ & 变色龙座（Chamaeleon） & 2000 年发现盘；首个确认的棕矮星盘，同时也是首个发射 X 射线的棕矮星 $^{[164]}$\\
\hline
首个 T 型双星棕矮星 & Epsilon Indi Ba、Bb$^{[165]}$T1 + T6 & 22h03m21.65363s，−56°47′09.5228″ & 印第安座（Indus） & 距离：3.626 pc\\
\hline
首个银晕棕矮星 & 2MASS J05325346+8246465 & sdL7 & 05h32m53.46s，+82°46′46.5″ & 双子座（Gemini） & Burgasser 等，2003 年 $^{[166]}$\\
\hline
首个伴随行星质量天体（planemo）的棕矮星 & 2M1207 & M8 & 12h07m33.47s，−39°32′54.0″ & 半人马座（Centaurus）& 2004 年发现 planemo，2005 年确认\\
\hline
首个具有双极外流的棕矮星 & Rho-Oph 102（SIMBAD：[GY92] 102） &  & 16h26m42.758s，−24°41′22.24″ & 蛇夫座（Ophiuchus）& 外流部分分辨 $^{[167]}$\\
\hline
首个食双星棕矮星系统 & 2M0535−05$^{[115][116]}$ & M6.5 & 15h35m21.84732s，−05°46′08.5714″ & 猎户座（Orion）& Stassun，2006、2007 年（距离约 450 pc）\\
\hline
首个被确认在主星红巨星阶段后仍幸存的棕矮星 & WD 0137−349 B $^{[168]}$ & L8 & 01h39m42.847s，−34°42′39.32″ & 雕刻室座（Sculptor） & 2006 年\\
\hline
最晚期 T 型光谱 & ULAS J003402.77−005206.7 & T9 $^{[69]}$ &  & 鲸鱼座（Cetus） & 2007 年\\
\hline
首个 Y 型光谱棕矮星 & CFBDS0059$^{[68]}$ & ≈Y0 & 00h59m10.83s，−01°14′01.3″ & 鲸鱼座（Cetus） & 2008 年；由于与其他 T 型棕矮星极为相似，也被归类为 T9 $^{[69]}$\\
\hline
首次探测到棕矮星差异自转 & TVLM 513−46546 & M9 & 15h01m08.3s，+22°50′02″ & 牧夫座（Boötes）& 赤道自转速度比两极快 0.022 弧度/天 $^{[169]}$\\
\hline
首次发现潜在棕矮星极光 & LSR J1835+3259 & M8.5 &  & 天琴座（Lyra） & 2015 年\\
\hline
首个具有大尺度赫比格–哈罗天体的棕矮星 & Mayrit 1701117（赫比格–哈罗天体：HH 1165）& 原棕矮星（proto-BD） & 05h40m25.799s，−02°48′55.42″ & 猎户座（Orion） & 赫比格–哈罗天体投影长度约 0.8 光年（0.26 pc）$^{[137]}$\\
\hline
首个 Y 型双星棕矮星系统 & WISE J0336−0143 & Y + Y & 03h36m05.052s，−01°43′50.48″ & 波江座（Eridanus）& 2023 年 $^{[108]}$\\
\hline
\end{tabular}
\end{table}
\subsubsection{极值一览表}
\begin{table}[ht]
\centering
\caption{请输入表格标题}\label{tab_Brdw2}
\begin{tabular}{|c|c|c|c|c|c|}
\hline
\textbf{记录} & \textbf{名称} & \textbf{光谱型} & \textbf{赤经/赤纬} & \textbf{星座} &\textbf{备注} \\
\hline
Oldest（最古老） & TOI-7019 b  &  &  18h 15m 09.436s，+50° 51′ 46.98″$^{[170]}$ & 天龙座（Draco）$^{[171]}$ & 年龄为 120±20 亿年$^{[172]}$\\
\hline
Youngest（最年轻） & 2MASS J05413280-0151272  & M8.5  & 05h 41m 32.801s，−01° 51′ 27.20″ & 猎户座（Orion） & 约 50 万年历史的火焰星云（Flame Nebula）成员之一的棕矮星，质量约 20.9 个木星质量（MJ）$^{[173]}$\\
\hline
Most massive（质量最大） & SDSS J010448.46+153501.8 $^{[174]}$ &usdL1.5 & 01h 04m 48.46s，+15° 35′ 01.8″  & 双鱼座（Pisces） & 距离约 180–290 秒差距（pc），质量约 88.5–91.7 个木星质量（MJ），属于过渡型棕矮星\\
\hline
Metal-rich（富金属） &   &   &   &   & \\
\hline
Metal-poor（贫金属） & TOI-7019 b  & 18h 15m 09.436s，+50° 51′ 46.98″$^{[170]}$& 天龙座（Draco）$^{[171]}$ & 金属丰度$[{\rm Fe/H}] = -0.79 \pm 0.05$［172］\\
\hline
Least massive（质量最小）&   &   &   &   & \\
\hline
Largest（半径最大） & FU Tauri A  & M7.25 & 04h 23m 35s，+25° 03′ 02″  & 金牛座（Taurus） & 半径为 1.803 $R_\odot$（约 1,254,000 km）$^{[175]}$\\
\hline
mallest（半径最小） & ZTF J1406+1222 B &  & 14h 06m 56s，−12° 22′ 43″  & 牧夫座（Boötes）$^{[176]}$ & 半径为 0.029 $R_\odot$（约 20,200 km）$^{[177]}$\\
\hline
Fastest rotating（自转最快） & 2MASS J03480772−6022270  & T7 &03h 48m 07.72s，−60° 22′ 27.1″  & 网罟座（Reticulum） & 自转周期 $1.080(^{+0.004}_{-0.005})$ 小时［178］\\
\hline
Farthest（距离最远） & 小麦哲伦云中的棕矮星候选体 &   & 01h 29m 32s，−73° 33′ 38″［a］ & 水蛇座（Hydrus） & 距离约 199,000 光年$^{[179]}$\\
\hline
Nearest（距离最近） & Luhman 16 AB  & L7.5 + T0.5 ± 1 & 10h 49m 18.723s，−53° 19′ 09.86″ & 船帆座（Vela） & 距离约 6.5 光年\\
\hline
Apparently brightest（表观最亮） & LP 944-20 & 光学：M9β；红外：L0 & 03h 39m 35.220^s，−35° 25′ 44.09″& 天炉座（Fornax） & 根据“超冷天体基本物理性质”$^{[180]}$，该天体显示出年轻特征，因此可能是一颗质量为 19.85 ± 13.02 个木星质量（MJ）的棕矮星，其 $J_{\rm MKO} = 10.68 \pm 0.03$等\\

\hline
* & * & * & * & * & * \\
\hline
* & * & * & * & * & * \\
\hline
\end{tabular}
\end{table}
\subsection{另见}
\begin{itemize}
\item Fusor（天文学）
\item 棕矮星荒漠（Brown-dwarf desert）——天文学中的一个概念
\item 热木星（Hot Jupiter）
\item 蓝矮星（红矮星阶段）——由红矮星演化而来的假想恒星类型
\item 暗物质——假想的不可见宇宙物质
\item 系外行星——太阳系之外的行星
\item 恒星化（Stellification）
\item 亚棕矮星（Sub-brown dwarf）
\item WD 0032−317 b
\item 棕矮星列表
\item Y 型棕矮星列表
\end{itemize}
\subsection{脚注}
a.发现这些棕矮星的疏散星团 NGC 602的位置
\subsection{参考文献}
\begin{enumerate}
\item Sorahana, Satoko；Yamamura, Issei；Murakami, Hiroshi（2013）。“基于 AKARI 近红外光谱测量的棕矮星半径研究”。《天体物理学杂志》（The Astrophysical Journal），767（1）：77。arXiv:1304.1259。书目代码：2013ApJ...767...77S。doi:10.1088/0004-637X/767/1/77。我们发现棕矮星的半径范围为 0.64–1.13 (R_J)，平均半径为 0.83 (R_J)。
\item Boss, Alan；McDowell, Tina（2001 年 4 月 3 日）。“它们是行星，还是别的什么？”未定名文档。华盛顿卡内基研究所。原文档于 2006 年 9 月 28 日存档。检索日期：2022 年 3 月 31 日。
\item Burgasser, Adam J.（2008 年 6 月）。“棕矮星：失败的恒星，或超级木星”（PDF）。《今日物理》（Physics Today），61（6）。马萨诸塞州剑桥：麻省理工学院出版社，第 70–71 页。书目代码：2008PhT....61f..70B。doi:10.1063/1.2947658。PDF 原文于 2013 年 5 月 8 日存档。检索日期：2022 年 3 月 31 日——经由美国物理学会（American Institute of Physics）。
\item Springer, Cham（2014）。Joergens, Viki（编）。《棕矮星五十年》（50 Years of Brown Dwarfs）。《天体物理与空间科学文库》（Astrophysics and Space Science Library），第 401 卷。SpringerLink。前言 xi 页，正文 168 页。doi:10.1007/978-3-319-01162-2。eISSN 2214-7985。ISBN 978-3-319-01162-2。ISSN 0067-0057。检索日期：2022 年 3 月 31 日。
\item Cain, Fraser（2009 年 1 月 6 日）。“如果‘棕色’不是一种颜色，那棕矮星是什么颜色？”检索日期：2013 年 9 月 24 日。
\item Burrows, Adam；Hubbard, William B.；Lunine, Jonathan I.；Liebert, James（2001）。“棕矮星与系外巨行星理论”。《现代物理评论》（Reviews of Modern Physics），73（3）：719–765。arXiv:astro-ph/0103383。书目代码：2001RvMP...73..719B。doi:10.1103/RevModPhys.73.719。S2CID 204927572。
\item O’Neill, Ian（2011 年 9 月 13 日）。“附近棕矮星上肆虐的剧烈风暴”。Seeker.com。
\item Finley, Dave。“研究表明：棕矮星与恒星共享相同的形成过程”。美国国家射电天文台（The National Radio Astronomy Observatory）。检索日期：2025 年 2 月 11 日。
\item Kumar, Shiv S.（1962）。“极低质量恒星中简并性的研究”。《天文学杂志》（Astronomical Journal），67：579。书目代码：1962AJ.....67S.579K。doi:10.1086/108658。
\item Tarter, Jill（2014）。“棕色并非一种颜色：‘棕矮星’一词的引入”。载于 Joergens, Viki（编），《棕矮星五十年：从预测到发现再到研究前沿》。 《天体物理与空间科学文库》，第 401 卷。Springer，第 19–24 页。doi:10.1007/978-3-319-01162-2_3。ISBN 978-3-319-01162-2。
\item Croswell, Ken（1999）。《行星探索：外星太阳系的史诗式发现》（Planet Quest: The Epic Discovery of Alien Solar Systems）。牛津大学出版社，第 118–119 页。ISBN 978-0-192-88083-3。
\item “太阳何时会变成黑矮星？”《天文学杂志》（Astronomy Magazine）。2020 年 4 月 10 日。检索日期：2022 年 5 月 2 日——经由 Astronomy.com。
\item Kumar, Shiv S.（1963）。“极低质量恒星的结构”。《天体物理学杂志》（Astrophysical Journal），137：1121。书目代码：1963ApJ...137.1121K。doi:10.1086/147589。
\item Burgasser, Adam J.（2008 年 6 月 1 日）。“棕矮星：失败的恒星，或超级木星”。《今日物理》（Physics Today），61（6）：70–71。书目代码：2008PhT....61f..70B。doi:10.1063/1.2947658。ISSN 003
\item Hayashi, Chushiro；Nakano, Takenori（1963）。“主序前阶段低质量恒星的演化”。《理论物理进展》（Progress of Theoretical Physics），30（4）：460–474。书目代码：1963PThPh..30..460H。doi:10.1143/PTP.30.460。
\item Nakano, Takenori（2014）。“主序前演化与氢燃烧的最小质量”。载于 Joergens, Viki（编），《棕矮星五十年：从预测到发现再到研究前沿》。 《天体物理与空间科学文库》，第 401 卷。Springer，第 5–17 页。doi:10.1007/978-3-319-01162-2_2。ISBN 978-3-319-01162-2。S2CID 73521636。
\item Martín, Eduardo L.；Basri, Gibor；Delfosse, Xavier；Forveille, Thierry（1997）。“棕矮星 DENIS-P J1228.2-1547 的 Keck HIRES 光谱”。《天文学与天体物理学》（Astronomy and Astrophysics），327：L29–L32。书目代码：1997A&A...327L..29M。
\item Kirkpatrick, J. Davy；Reid, I. Neill；Liebert, James；Cutri, Roc M.；Nelson, Brant；Beichmann, Charles A.；Dahn, Conard C.；Monet, David G.；Gizis, John E.；Skrutskie, Michael F.（1999）。“比 M 型更冷的矮星：基于二维微米全天巡天（2MASS）发现对 L 型光谱类型的定义”（PDF）。《天体物理学杂志》（The Astrophysical Journal），519（2）：802–833。书目代码：1999ApJ...519..802K。doi:10.1086/307414。S2CID 73569208。
\item Zuckerman, B.；Becklin, E. E.（1992 年 2 月）。“白矮星的伴星——极低质量恒星与棕矮星候选体 GD 165B”。《天体物理学杂志》（The Astrophysical Journal），386：260。书目代码：1992ApJ...386..260Z。doi:10.1086/171012。ISSN 0004-637X。
\item Kirkpatrick, J. D.；Henry, Todd J.；Liebert, James（1993 年 4 月）。“棕矮星候选体 GD 165B 的独特光谱及其与其他低光度天体光谱的比较”。《天体物理学杂志》（The Astrophysical Journal），406：701。书目代码：1993ApJ...406..701K。doi:10.1086/172480。ISSN 0004-637X。
\item “加那利群岛天体物理研究所（Instituto de Astrofísica de Canarias, IAC）”。Iac.es。检索日期：2013-03-16。
\item Rebolo, Rafael（2014）。“Teide 1 与棕矮星的发现”。载于 Joergens, Viki（编），《棕矮星五十年——从预测到发现再到研究前沿》。天体物理与空间科学文库，第 401 卷。Springer，第 25–50 页。书目代码：2014ASSL..401...25R。doi:10.1007/978-3-319-01162-2_4。ISBN 978-3-319-01162-2。
\item Rebolo, Rafael；Zapatero-Osorio, María Rosa；Martín, Eduardo L.（1995 年 9 月）。“在昴宿星团中发现一颗棕矮星”。《自然》（Nature），377（6545）：129–131。书目代码：1995Natur.377..129R。doi:10.1038/377129a0。S2CID 28029538。
\item Leech, Kieron；Altieri, Bruno；Metcalfe, Liam；Martín, Eduardo L.；Rebolo, Rafael；Zapatero-Osorio, María Rosa；Laureijs, René J.；Prusti, Timo；Salama, Alberto；Siebenmorgen, Ralf；Claes, Peter；Trams, Norman（2000）。“昴宿星团棕矮星 Teide 1 与 Calar 3 的中红外观测”。ASP 会议论文集（ASP Conference Series），212：82–87。书目代码：2000ASPC..212...82L。
\item Rebolo, R.；Martín, E. L.；Basri, G.；Marcy, G. W.；Zapatero-Osorio, M. R.（1996）。“通过锂检验确认的昴宿星团棕矮星”。《天体物理学杂志》（The Astrophysical Journal），469：53–56。arXiv:astro-ph/9607002。书目代码：1996ApJ...469L..53R。doi:10.1086/310263。S2CID 119485127。
\item Kirkpatrick, J. Davy；Burgasser, Adam J.（2012 年 11 月 6 日）。“M、L 与 T 型矮星的测光、光谱与天体测量”。DwarfArchives.org。加利福尼亚州帕萨迪纳：加州理工学院。检索日期：2012-12-28。（M=536，L=918，T=355，Y=14）
\item McCaughrean, Mark J.；Close, Laird M.；Scholz, Ralf-Dieter；Lenzen, Rainer；Biller, Beth A.；Brandner, Wolfgang；Hartung, Markus；Lodieu, Nicolas（2004 年 1 月）。“Epsilon Indi Ba/Bb：最近的双星棕矮星系统”。《天文学与天体物理学》（Astronomy & Astrophysics），413（3）：1029–1036。arXiv:astro-ph/0309256。doi:10.1051/0004-6361:20034292。S2CID 15407249。
\item “天文学家宣布发现棕矮星的首个明确证据”。太空望远镜科学研究所（STScI）。检索日期：2019-10-23。
\item Forbes, John C.；Loeb, Abraham（2019 年 2 月）。“关于质量超过氢燃烧极限的棕矮星是否存在”。《天体物理学杂志》（The Astrophysical Journal），871（2）：11。arXiv:1805.12143。书目代码：2019ApJ...871..227F。doi:10.3847/1538-4357/aafac8。S2CID 119059288，第 227 页。
\item Burrows, Adam；Hubbard, W. B.；Lunine, J. I.；Liebert, James（2001 年 7 月）。“棕矮星与系外巨行星的理论”。《现代物理评论》（Reviews of Modern Physics），73（3）：719–765。arXiv:astro-ph/0103383。书目代码：2001RvMP...73..719B。doi:10.1103/RevModPhys.73.719。S2CID 204927572。由此，在太阳金属丰度且 $Y_\alpha = 50.25$ 条件下，氢燃烧最小质量（HBMM）为 0.07–0.074 $M_\odot$，而在零金属丰度情况下为 0.092$M_\odot$。
\item Martín, E. L.；Lodieu, N.；del Burgo, C.（2022-02-21）。“对棕矮星中热核锂燃烧最小质量的新约束”。《皇家天文学会月刊》（Monthly Notices of the Royal Astronomical Society），510（2）：2841–2850。doi:10.1093/mnras/stab2969。ISSN 0035-8711。
\item Kulkarni, Shrinivas R.（1997 年 5 月 30 日）。“棕矮星：可能连接恒星与行星的缺失环节”。《科学》（Science），276（5317）：1350–1354。书目代码：1997Sci...276.1350K。doi:10.1126/science.276.5317.1350。
\item Burgasser, Adam J.；Marley, Mark S.；Ackerman, Andrew S.；Saumon, Didier；Lodders, Katharina；Dahn, Conard C.；Harris, Hugh C.；Kirkpatrick, J. Davy（2002-06-01）。“L/T 型矮星过渡区云层破碎的证据”。《天体物理学杂志》（The Astrophysical Journal），571（2）：L151–L154。arXiv:astro-ph/0205051。书目代码：2002ApJ...571L.151B。doi:10.1086/341343。ISSN 0004-637X。
\item Vos, Johanna M.；Faherty, Jacqueline K.；Gagné, Jonathan；Marley, Mark；Metchev, Stanimir；Gizis, John；Rice, Emily L.；Cruz, Kelle（2022-01-01）。“让伟大的世界旋转：利用斯皮策空间望远镜揭示年轻巨行星类比体的风暴性与湍动特征”。《天体物理学杂志》（The Astrophysical Journal），924（2）：68。arXiv:2201.04711。书目代码：2022ApJ...924...68V。doi:10.3847/1538-4357/ac4502。ISSN 0004-637X。
\item Vos, Johanna M.；Burningham, Ben；Faherty, Jacqueline K.；Alejandro, Sherelyn；Gonzales, Eileen；Calamari, Emily；Bardalez Gagliuffi, Daniella；Visscher, Channon；Tan, Xianyu；Morley, Caroline V.；Marley, Mark；Gemma, Marina E.；Whiteford, Niall；Gaarn, Josefine；Park, Grace（2023-02-01）。“两颗高度变异的系外行星类比体大气中的斑块状镁橄榄石云”。《天体物理学杂志》（The Astrophysical Journal），944（2）：138。arXiv:2212.07399。书目代码：2023ApJ...944..138V。doi:10.3847/1538-4357/acab58。ISSN 0004-637X。
\item Manjavacas, Elena；Karalidi, Theodora；Vos, Johanna M.；Biller, Beth A.；Lew, Ben W. P.（2021-11-01）。“通过 Keck I/MOSFIRE 光谱测光变率揭示一颗年轻低质量棕矮星的垂直云结构——其为直接成像系外行星 β-绘架座 b 的类比体”。《天文学杂志》（The Astronomical Journal），162（5）：179。arXiv:2107.12368。书目代码：2021AJ....162..179M。doi:10.3847/1538-3881/ac174c。ISSN 0004-6256。
\item Tremblin, P.；Chabrier, G.；Baraffe, I.；Liu, Michael C.；Magnier, E. A.；Lagage, P.-O.；Alves de Oliveira, C.；Burgasser, A. J.；Amundsen, D. S.；Drummond, B.（2017-11-01）。“年轻、低重力亚恒星天体的无云大气”。《天体物理学杂志》（The Astrophysical Journal），850（1）：46。arXiv:1710.02640。书目代码：2017ApJ...850...46T。doi:10.3847/1538-4357/aa9214。ISSN 0004-637X。
\item Suárez, Genaro；Metchev, Stanimir（2022-07-01）。“使用斯皮策红外光谱仪观测的超冷矮星——II：L 型矮星中硅酸盐云的出现与沉降，以及完整 M5–T9 场矮星光谱样本的分析”。《皇家天文学会月刊》（Monthly Notices of the Royal Astronomical Society），513（4）：5701–5726。arXiv:2205.00168。书目代码：2022MNRAS.513.5701S。doi:10.1093/mnras/stac1205。ISSN 0035-8711。
\item Miles, Brittany E.；Biller, Beth A.；Patapis, Polychronis；Worthen, Kadin；Rickman, Emily；Hoch, Kielan K. W.；Skemer, Andrew；Perrin, Marshall D.；Whiteford, Niall；Chen, Christine H.；Sargent, B.；Mukherjee, Sagnick；Morley, Caroline V.；Moran, Sarah E.；Bonnefoy, Mickael（2023-03-01）。“JWST 系外行星系统直接观测早期科学计划 II：行星质量伴星 VHS 1256-1257 b 在 1–20 μm 波段的光谱”。《天体物理学杂志》（The Astrophysical Journal），946（1）：L6。arXiv:2209.00620。书目代码：2023ApJ...946L...6M。doi:10.3847/2041-8213/acb04a。ISSN 0004-637X。
\item Biller, Beth A.；Crossfield, Ian J. M.；Mancini, Luigi；Ciceri, Simona；Southworth, John；Kopytova, Taisiya G.；Bonnefoy, Mickaël；Deacon, Niall R.；Schlieder, Joshua E.；Buenzli, Esther；Brandner, Wolfgang；Allard, France；Homeier, Derek；Freytag, Bernd；Bailer-Jones, Coryn A. L.；Greiner, Jochen；Henning, Thomas；Goldman, Bertrand（2013 年 11 月 6 日）。“最近棕矮星的天气：对 WISE J104915.57–531906.1AB 的分辨式同步多波段变率监测”。《天体物理学杂志快报》（The Astrophysical Journal Letters），778（1）：L10。arXiv:1310.5144。书目代码：2013ApJ...778L..10B。doi:10.1088/2041-8205/778/1/L10。S2CID 56107487。
\item Morley, Caroline V.；Fortney, Jonathan J.；Marley, Mark S.；Visscher, Channon；Saumon, Didier；Leggett, S. K.（2012-09-01）。“T 型与 Y 型矮星大气中被忽视的云层”。《天体物理学杂志》（The Astrophysical Journal），756（2）：172。arXiv:1206.4313。书目代码：2012ApJ...756..172M。doi:10.1088/0004-637X/756/2/172。ISSN 0004-637X。S2CID 118398946。
\item Manjavacas, Elena；Karalidi, Theodora；Tan, Xianyu；Vos, Johanna M.；Lew, Ben W. P.；Biller, Beth A.；Oliveros-Gómez, Natalia（2022-08-01）。“快速自转的晚期 T 型矮星：大气顶层与垂直云结构”。《天文学杂志》（The Astronomical Journal），164（2）：65。arXiv:2206.07566。书目代码：2022AJ....164...65M。doi:10.3847/1538-3881/ac7953。ISSN 0004-6256。
\item Faherty, Jacqueline K.；Tinney, C. G.；Skemer, Andrew；Monson, Andrew J.（2014-09-01）。“在已知最冷的棕矮星中发现水云存在的迹象”。《天体物理学杂志》（The Astrophysical Journal），793（1）：L16。arXiv:1408.4671。书目代码：2014ApJ...793L..16F。doi:10.1088/2041-8205/793/1/L16。hdl:1959.4/unsworks_36908。ISSN 0004-637X。S2CID 119246100。
\item Basri, Gibor；Brown, Michael E.（2006-08-20）。“从微行星到棕矮星：什么是行星？”《地球与行星科学年评》（Annual Review of Earth and Planetary Sciences），34（2006）：193–216。arXiv:astro-ph/0608417。书目代码：2006AREPS..34..193B。doi:10.1146/annurev.earth.34.031405.125058。S2CID 119338327。
\item Chen, Jingjing；Kipping, David（2016）。“其他世界质量与半径的概率预测”。《天体物理学杂志》（The Astrophysical Journal），834（1）：17。arXiv:1603.08614。书目代码：2017ApJ...834...17C。doi:10.3847/1538-4357/834/1/17。S2CID 119114880。
\item 《类木行星：天王星与海王星》。原文档于 2012-01-18 存档。检索日期：2013-03-15。
\item 《Cool Cosmos——行星与卫星》。原文档于 2019-02-21 存档。检索日期：2019-02-11。
\item “系外行星工作组：‘行星’的定义”。国际天文学联合会（IAU）立场声明。2003-02-28。原文档于 2014-12-16 存档。检索日期：2014-04-28。
\item Bodenheimer, Peter；D’Angelo, Gennaro；Lissauer, Jack J.；Fortney, Jonathan J.；Saumon, Didier（2013）。“由核吸积形成的高质量巨行星与低质量棕矮星中的氘燃烧”。《天体物理学杂志》（The Astrophysical Journal），770（2）：120（13 页）。arXiv:1305.0980。书目代码：2013ApJ...770..120B。doi:10.1088/0004-637X/770/2/120。S2CID 118553341。
\item Spiegel, David S.；Burrows, Adam；Milson, John A.（2011）。“棕矮星与巨行星的氘燃烧质量下限”。《天体物理学杂志》（The Astrophysical Journal），727（1）：57。arXiv:1008.5150。书目代码：2011ApJ...727...57S。doi:10.1088/0004-637X/727/1/57。S2CID 118513110。
\item Schneider, Jean；Dedieu, Cyril；Le Sidaner, Pierre；Savalle, Renaud；Zolotukhin, Ivan（2011）。“系外行星的定义与编目：exoplanet.eu 数据库”。《天文学与天体物理学》（Astronomy & Astrophysics），532（79）：A79。arXiv:1106.0586。书目代码：2011A&A...532A..79S。doi:10.1051/0004-6361/201116713。S2CID 55994657。
\item Schneider, Jean（2016 年 7 月）。“系外行星与棕矮星：CoRoT 的视角与未来”。《CoRoT 遗产之书》（The CoRoT Legacy Book），第 157 页。arXiv:1604.00917。doi:10.1051/978-2-7598-1876-1.c038。ISBN 978-2-7598-1876-1。S2CID 118434022。
\item Hatzes, Artie P.；Rauer, Heike（2015）。“基于质量—密度关系的巨行星定义”。《天体物理学杂志》（The Astrophysical Journal），810（2）：L25。arXiv:1506.05097。书目代码：2015ApJ...810L..25H。doi:10.1088/2041-8205/810/2/L25。S2CID 119111221。
\item Wright, Jason T.；Fakhouri, Onsi；Marcy, Geoffrey W.；Han, Eunkyu；Feng, Y. Katherina；Johnson, John Asher；Howard, Andrew W.；Fischer, Debra A.；Valenti, Jeff A.；Anderson, Jay；Piskunov, Nikolai（2010）。“系外行星轨道数据库”。《太平洋天文学会出版物》（Publications of the Astronomical Society of the Pacific），123（902）：412–422。arXiv:1012.5676。书目代码：2011PASP..123..412W。doi:10.1086/659427。S2CID 51769219。
\item 《NASA 系外行星档案：纳入档案的系外行星标准》。
\item “系外行星工作组——‘行星’的定义”。IAU 关于“行星”定义的立场声明。原文档于 2012-07-02 存档于 Wayback Machine。
\item Delorme, Philippe；Gagné, Jonathan；Malo, Lison；Reylé, Céline；Artigau, Étienne；Albert, Loïc；Forveille, Thierry；Delfosse, Xavier；Allard, France；Homeier, Derek（2012 年 12 月）。“CFBDSIR2149-0403：AB 多拉多斯年轻移动群中的一颗 4–7 木星质量的自由漂浮行星？”《天文学与天体物理学》（Astronomy & Astrophysics），548：A26。arXiv:1210.0305。书目代码：2012A&A...548A..26D。doi:10.1051/0004-6361/201219984。S2CID 50935950。
\item Luhman, Kevin L.（2014 年 4 月 21 日）。“在距离太阳 2 pc 处发现一颗约 250 K 的棕矮星”。《天体物理学杂志快报》（The Astrophysical Journal Letters），786（2）：L18。arXiv:1404.6501。书目代码：2014ApJ...786L..18L。doi:10.1088/2041-8205/786/2/L18。S2CID 119102654。
\item Saumon, Didier；Marley, Mark S.（2008 年 12 月）。“L 型与 T 型矮星在颜色—星等图中的演化”。《天体物理学杂志》（Astrophysical Journal），689（2）：1327–1344。arXiv:0808.2611。书目代码：2008ApJ...689.1327S。doi:10.1086/592734。ISSN 0004-637X。S2CID 15981010。
\item Marocco, Federico；Kirkpatrick, J. Davy；Meisner, Aaron M.；Caselden, Dan；Eisenhardt, Peter R. M.；Cushing, Michael C.；Faherty, Jacqueline K.；Gelino, Christopher R.；Wright, Edward L.（2020）。“对极端寒冷棕矮星 CWISEP J144606.62−231717.8 的改进红外测光与初步视差测量”。《天体物理学杂志》（The Astrophysical Journal），888（2）：L19。arXiv:1912.07692。书目代码：2020ApJ...888L..19M。doi:10.3847/2041-8213/ab6201。S2CID 209386563。
\item Filippazzo, Joseph C.；Rice, Emily L.；Faherty, Jacqueline K.；Cruz, Kelle L.；Van Gordon, Mollie M.；Looper, Dagny L.（2015 年 9 月）。“覆盖恒星到行星质量范围的年轻与场星年龄天体：基本参数与光谱能量分布”。《天体物理学杂志》（Astrophysical Journal），810（2）：158。arXiv:1508.01767。书目代码：2015ApJ...810..158F。doi:10.1088/0004-637X/810/2/158。ISSN 0004-637X。S2CID 89611607。
\item Mohanty, Subhanjoy；Jayawardhana, Ray；Huélamo, Nuria；Mamajek, Eric（2007 年 3 月）。“行星质量伴星 2MASS 1207-3932B：温度、质量及边缘朝向盘面的证据”。《天体物理学杂志》（Astrophysical Journal），657（2）：1064–1091。arXiv:astro-ph/0610550。书目代码：2007ApJ...657.1064M。doi:10.1086/510877。ISSN 0004-637X。S2CID 17326111。
\item Zhang, Zhoujian；Liu, Michael C.；Hermes, James J.；Magnier, Eugene A.；Marley, Mark S.；Tremblay, Pier-Emmanuel；Tucker, Michael A.；Do, Aaron；Payne, Anna V.；Shappee, Benjamin J.（2020 年 2 月）。“超宽轨道上的冷伴星（COCONUTS）I：一颗老年白矮星周围的高重力 T4 基准体，以及对 L/T 过渡的表面重力依赖性的再检视”。《天体物理学杂志》（The Astrophysical Journal），891（2）：171。arXiv:2002.05723。书目代码：2020ApJ...891..171Z。doi:10.3847/1538-4357/ab765c。S2CID 211126544。
\item Smart, Richard L.；Bucciarelli, Beatrice；Jones, Hugh R. A.；Marocco, Federico；Andrei, Alexandre Humberto；Goldman, Bertrand；Méndez, René A.；d’Avila, Victor de A.；Burningham, Ben；Camargo, Julio Ignácio Bueno de；Crosta, Maria Teresa；Daprà, Mario；Jenkins, James S.；Lachaume, Regis；Lattanzi, Mario G.；Penna, Jucira L.；Pinfield, David J.；da Silva Neto, Dario Nepomuceno；Sozzetti, Alessandro；Vecchiato, Alberto（2018 年 12 月）。“南天极端冷天体视差 III：118 颗 L 型与 T 型矮星”。《皇家天文学会月刊》（MNRAS），481（3）：3548–3562。arXiv:1811.00672。书目代码：2018MNRAS.481.3548S。doi:10.1093/mnras/sty2520。ISSN 0035-8711。S2CID 119390019。
\item Burrows, Adam；Hubbard, William B.；Lunine, Jonathan I.；Liebert, James（2001）。“棕矮星与系外巨行星的理论”。《现代物理评论》（Reviews of Modern Physics），73（3）：719–765。arXiv:astro-ph/0103383。书目代码：2001RvMP...73..719B。doi:10.1103/RevModPhys.73.719。S2CID 204927572。
\item “棕矮星类型的艺术家视图”。原文档于 2011-11-17 存档于 Wayback Machine。
\item Leggett, Sandy K.；Cushing, Michael C.；Saumon, Didier；Marley, Mark S.；Roellig, Thomas L.；Warren, Stephen J.；Burningham, Ben；Jones, Hugh R. A.；Kirkpatrick, J. Davy；Lodieu, Nicolas；Lucas, Philip W.；Mainzer, Amy K.；Martín, Eduardo L.；McCaughrean, Mark J.；Pinfield, David J.；Sloan, Gregory C.；Smart, Richard L.；Tamura, Motohide；Van Cleve, Jeffrey（2009）。“四颗约 600 K 的 T 型矮星的物理性质”。《天体物理学杂志》（The Astrophysical Journal），695（2）：1517–1526。arXiv:0901.4093。书目代码：2009ApJ...695.1517L。doi:10.1088/0004-637X/695/2/1517。S2CID 44050900。
\item Delorme, Philippe；Delfosse, Xavier；Albert, Loïc；Artigau, Étienne；Forveille, Thierry；Reylé, Céline；Allard, France；Homeier, Derek；Robin, Annie C.；Willott, Chris J.；Liu, Michael C.；Dupuy, Trent J.（2008）。“CFBDS J005910.90-011401.3：触及 T–Y 棕矮星过渡区？”《天文学与天体物理学》（Astronomy and Astrophysics），482（3）：961–971。arXiv:0802.4387。书目代码：2008A&A...482..961D。doi:10.1051/0004-6361:20079317。S2CID 847552。
\item Burningham, Ben；Pinfield, David J.；Leggett, Sandy K.；Tamura, Motohide；Lucas, Philip W.；Homeier, Derek；Day-Jones, Avril；Jones, Hugh R. A.；Clarke, J. R. A.；Ishii, Miki；Kuzuhara, Masayuki；Lodieu, Nicolas；Zapatero-Osorio, María Rosa；Venemans, Bram Pieter；Mortlock, Daniel J.；Barrado y Navascués, David；Martín, Eduardo L.；Magazzù, Antonio（2008）。“探索亚恒星温度区间，最低至约 550 K”。《皇家天文学会月刊》（Monthly Notices of the Royal Astronomical Society），391（1）：320–333。arXiv:0806.0067。书目代码：2008MNRAS.391..320B。doi:10.1111/j.1365-2966.2008.13885.x。S2CID 1438322。
\item Beiler, Samuel A.；Cushing, Michael C.；Kirkpatrick, J. Davy；Schneider, Adam C.；Mukherjee, Sagnick；Marley, Mark S.（2023-07-01）。“JWST 首次给出一颗 Y 型矮星的光谱能量分布”。《天体物理学杂志》（The Astrophysical Journal），951（2）：L48。arXiv:2306.11807。书目代码：2023ApJ...951L..48B。doi:10.3847/2041-8213/ace32c。ISSN 0004-637X。
\item Luhman, K. L.；Alves de Oliveira, C.（2025 年 6 月）。“在 IC 348 中、位于初始质量函数底端的棕矮星新光谱类别”。《天体物理学杂志》（The Astrophysical Journal），986（1）：L14。arXiv:2506.08969。书目代码：2025ApJ...986L..14L。doi:10.3847/2041-8213/addc55。ISSN 0004-637X。
\item Zahnle, Kevin J.；Marley, Mark S.（2014-12-01）。“棕矮星与自发光巨行星中的甲烷、一氧化碳与氨”。《天体物理学杂志》（The Astrophysical Journal），797（1）：41。arXiv:1408.6283。书目代码：2014ApJ...797...41Z。doi:10.1088/0004-637X/797/1/41。ISSN 0004-637X。S2CID 118509317。
\item Bardalez Gagliuffi, Daniella C.；Faherty, Jacqueline K.；Schneider, Adam C.；Meisner, Aaron；Caselden, Dan；Colin, Guillaume；Goodman, Sam；Kirkpatrick, J. Davy；Kuchner, Marc；Gagné, Jonathan；Logsdon, Sarah E.；Burgasser, Adam J.；Allers, Katelyn；Debes, John；Wisniewski, John（2020-06-01）。“WISEA J083011.95+283716.0：一枚缺失环节的行星质量天体”。《天体物理学杂志》（The Astrophysical Journal），895（2）：145。arXiv:2004.12829。书目代码：2020ApJ...895..145B。doi:10.3847/1538-4357/ab8d25。ISSN 0004-637X。S2CID 216553879。
\item “光谱类型代码”。simbad.u-strasbg.fr。检索日期：2020-03-06。
\item Burningham, Ben；Smith, Leigh；Cardoso, Cátia V.；Lucas, Philip W.；Burgasser, Adam J.；Jones, Hugh R. A.；Smart, Richard L.（2014 年 5 月）。“发现一颗 T6.5 次矮星”。《皇家天文学会月刊》（Monthly Notices of the Royal Astronomical Society），440（1）：359–364。arXiv:1401.5982。书目代码：2014MNRAS.440..359B。doi:10.1093/mnras/stu184。ISSN 0035-8711。S2CID 119283917。
\item Cruz, Kelle L.；Kirkpatrick, J. Davy；Burgasser, Adam J.（2009 年 2 月）。“在场星中识别出的年轻 L 型矮星：从 L0 到 L5 的初步低重力光学光谱序列”。《天文学杂志》（The Astronomical Journal），137（2）：3345–3357。arXiv:0812.0364。书目代码：2009AJ....137.3345C。doi:10.1088/0004-6256/137/2/3345。ISSN 0004-6256。S2CID 15376964。
\item Looper, Dagny L.；Kirkpatrick, J. Davy；Cutri, Roc M.；Barman, Travis；Burgasser, Adam J.；Cushing, Michael C.；Roellig, Thomas；McGovern, Mark R.；McLean, Ian S.；Rice, Emily；Swift, Brandon J.（2008 年 10 月）。“来自 2MASS 自行巡天的两颗邻近异常 L 型矮星的发现：年轻还是富金属？”《天体物理学杂志》（Astrophysical Journal），686（1）：528–541。arXiv:0806.1059。书目代码：2008ApJ...686..528L。doi:10.1086/591025。ISSN 0004-637X。S2CID 18381182。
\item Kirkpatrick, J. Davy；Looper, Dagny L.；Burgasser, Adam J.；Schurr, Steven D.；Cutri, Roc M.；Cushing, Michael C.；Cruz, Kelle L.；Sweet, Anne C.；Knapp, Gillian R.；Barman, Travis S.；Bochanski, John J.（2010 年 9 月）。“利用多历元二维微米全天巡天（2MASS）数据的近红外自行巡天发现”。《天体物理学杂志增刊》（Astrophysical Journal Supplement Series），190（1）：100–146。arXiv:1008.3591。书目代码：2010ApJS..190..100K。doi:10.1088/0067-0049/190/1/100。ISSN 0067-0049。S2CID 118435904。
\item Faherty, Jacqueline K.；Riedel, Adric R.；Cruz, Kelle L.；Gagne, Jonathan；Filippazzo, Joseph C.；Lambrides, Erini；Fica, Haley；Weinberger, Alycia；Thorstensen, John R.；Tinney, Chris G.；Baldassare, Vivienne（2016 年 7 月）。“作为系外行星类比体的棕矮星：族群性质”。《天体物理学杂志增刊》（Astrophysical Journal Supplement Series），225（1）：10。arXiv:1605.07927。书目代码：2016ApJS..225...10F。doi:10.3847/0067-0049/225/1/10。ISSN 0067-0049。S2CID 118446190。
\item Reid, Neill。“颜色—星等数据”。[www.stsci.edu。检索日期：2020-03-06。](http://www.stsci.edu。检索日期：2020-03-06。)
\item 美国国家射电天文台（National Radio Astronomy Observatory）（2020 年 4 月 9 日）。“天文学家测量棕矮星上的风速——其大气与内部以不同速度自转”。EurekAlert!。检索日期：2020 年 4 月 10 日。
\item Chen, Minghan；Li, Yiting；Brandt, Timothy D.；Dupuy, Trent J.；Cardoso, Cátia V.；McCaughrean, Mark J.（2022）。“ε Indi Ba 与 Bb 的精确动力学质量：L/T 过渡阶段冷却减缓的证据”。《天文学杂志》（The Astronomical Journal），163（6）：288。arXiv:2205.08077。书目代码：2022AJ....163..288C。doi:10.3847/1538-3881/ac66d2。S2CID 248834536。
\item “NASA 空间望远镜观测到棕矮星中的天气模式”。Hubblesite。NASA。原文档于 2014 年 4 月 2 日存档。检索日期：2013 年 1 月 8 日。
\item “天文学家测得太阳系外天体的强风速度”。CNN。2020 年 4 月 9 日。检索日期：2020 年 4 月 11 日。
\item Route, Matthew；Wolszczan, Alexander（2016 年 10 月 20 日）。“第二次使用阿雷西博望远镜搜索超冷矮星 5 GHz 射电耀斑”。《天体物理学杂志》（The Astrophysical Journal），830（2）：85。arXiv:1608.02480。书目代码：2016ApJ...830...85R。doi:10.3847/0004-637X/830/2/85。S2CID 119279978。
\item Rutledge, Robert E.；Basri, Gibor；Martín, Eduardo L.；Bildsten, Lars（2000 年 8 月 1 日）。“钱德拉探测到棕矮星 LP 944-20 的一次 X 射线耀斑”。《天体物理学杂志》（The Astrophysical Journal），538（2）：L141–L144。arXiv:astro-ph/0005559。书目代码：2000ApJ...538L.141R。doi:10.1086/312817。S2CID 17800872。
\item Berger, Edo；Ball, Steven；Becker, Kate M.；Clarke, Melanie；Frail, Dale A.；Fukuda, Therese A.；Hoffman, Ian M.；Mellon, Richard；Momjian, Emmanuel；Murphy, Nathanial W.；Teng, Stacey H.；Woodruff, Timothy；Zauderer, B. Ashley；Zavala, Robert T.（2001-03-15）。“发现棕矮星 LP944-20 的射电辐射”。《自然》（Nature，投稿手稿），410（6826）：338–340。arXiv:astro-ph/0102301。书目代码：2001Natur.410..338B。doi:10.1038/35066514。PMID 11268202。S2CID 4411256。原文档于 2021-04-27 存档。
\item Chauvin, Gael；Zuckerman, Ben；Lagrange, Anne-Marie。“是的，这确实是一张系外行星的图像：天文学家确认了太阳系外行星的首张影像”（新闻稿）。欧洲南方天文台（European Southern Observatory）。检索日期：2020-02-09。
\item Luhman, Kevin L.（2013 年 4 月）。“在距离太阳 2 pc 处发现一对棕矮星双星”。《天体物理学杂志快报》（Astrophysical Journal Letters），767（1）：L1。arXiv:1303.2401。书目代码：2013ApJ...767L...1L。doi:10.1088/2041-8205/767/1/L1。ISSN 0004-637X。S2CID 8419422。
\item Clavin, Whitney；Harrington, J. D.（2014 年 4 月 25 日）。“NASA 的斯皮策与 WISE 望远镜发现太阳附近一个距离很近、温度很低的邻居”。NASA.gov。原文档于 2014 年 4 月 26 日存档。
\item “来自棕矮星日冕的 X 射线”。2003 年 4 月 14 日。原文档于 2010 年 12 月 30 日存档。检索日期：2010 年 3 月 19 日。
\item Route, Matthew（2017 年 8 月 10 日）。“射电耀斑型超冷矮星族群综合研究”。《天体物理学杂志》（The Astrophysical Journal），845（1）：66。arXiv:1707.02212。书目代码：2017ApJ...845...66R。doi:10.3847/1538-4357/aa7ede。S2CID 118895524。
\item Kao, Melodie M.；Hallinan, Gregg；Pineda, J. Sebastian；Stevenson, David；Burgasser, Adam J.（2018 年 7 月 31 日）。“最冷的棕矮星上最强的磁场”。《天体物理学杂志增刊》（The Astrophysical Journal Supplement Series），237（2）：25。arXiv:1808.02485。书目代码：2018ApJS..237...25K。doi:10.3847/1538-4365/aac2d5。S2CID 118898602。
\item Route, Matthew（2017 年 7 月 10 日）。“WISEP J060738.65+242953.4 真的是一颗磁活动强、极向朝向我们的 L 型矮星吗？”《天体物理学杂志》（The Astrophysical Journal），843（2）：115。arXiv:1706.03010。书目代码：2017ApJ...843..115R。doi:10.3847/1538-4357/aa78ab。S2CID 119056418。
\item Route, Matthew（2016 年 10 月 20 日）。“在主序末端之外发现类似太阳的活动周期？”《天体物理学杂志快报》（The Astrophysical Journal Letters），830（2）：L27。arXiv:1609.07761。书目代码：2016ApJ...830L..27R。doi:10.3847/2041-8205/830/2/L27。S2CID 119111063。
\item Phys.org。“在一颗超冷恒星中发现破纪录的射电波”（新闻稿）。
\item Route, M.；Wolszczan, A.（2012 年 3 月 10 日）。“阿雷西博望远镜探测到最冷的射电耀斑棕矮星”。《天体物理学杂志快报》（The Astrophysical Journal Letters），747（2）：L22。arXiv:1202.1287。书目代码：2012ApJ...747L..22R。doi:10.1088/2041-8205/747/2/L22。S2CID 119290950。
\item Route, Matthew（2024 年 5 月 1 日）。“ROME IV：使用阿雷西博望远镜在 5 GHz 频段搜索被认为宿主系外行星系统中的亚恒星磁层射电辐射”。《天体物理学杂志》（The Astrophysical Journal），966（1）：55。arXiv:2403.02226。书目代码：2024ApJ...966...55R。doi:10.3847/1538-4357/ad30ff。
\item Meisner, Aaron；Kocz, Amanda。“绘制我们太阳邻域的地图”。NOIRLab。检索日期：2021 年 2 月 1 日。
\item O’Neill, Ian（2012 年 6 月 12 日）。“棕矮星——恒星‘垃圾堆’中的矮小成员，比想象中更为稀少”。Space.com。检索日期：2012-12-28。
\item Muzic, Koraljka；Schoedel, Rainer；Scholz, Alexander；Geers, Vincent C.；Jayawardhana, Ray；Ascenso, Joana；Cieza, Lucas A.（2017-07-02）。“大质量年轻星团 RCW 38 的低质量成员含量”。《皇家天文学会月刊》（Monthly Notices of the Royal Astronomical Society），471（3）：3699–3712。arXiv:1707.00277。书目代码：2017MNRAS.471.3699M。doi:10.1093/mnras/stx1906。ISSN 0035-8711。S2CID 54736762。
\item Apai, Dániel；Karalidi, T.；Marley, Mark S.；Yang, H.；Flateau, D.；Metchev, S.；Cowan, N. B.；Buenzli, E.；Burgasser, Adam J.；Radigan, J.；Artigau, Étienne；Lowrance, P.（2017）。“棕矮星大气中的区域、斑点以及行星尺度波动”。《科学》（Science），357（6352）：683–687。书目代码：2017Sci...357..683A。doi:10.1126/science.aam9848。PMID 28818943。
\item Gohd, Chelsea（2020 年 8 月 19 日）。“志愿者在太阳附近发现近 100 颗寒冷的棕矮星”。Space.com。
\item “外星天气报告：詹姆斯·韦布空间望远镜在两颗棕矮星上探测到炽热、含沙的强风”。Space.com。
\item “最冷的棕矮星是否多为孤立天体？”[www.noirlab.edu。检索日期：2023-04-16。](http://www.noirlab.edu。检索日期：2023-04-16。)
\item Fontanive, Clémence；Biller, Beth；Bonavita, Mariangela；Allers, Katelyn（2018-09-01）。“对最冷棕矮星多重性统计的约束：双星比例随光谱型继续下降”。《皇家天文学会月刊》（Monthly Notices of the Royal Astronomical Society），479（2）：2702–2727。arXiv:1806.08737。书目代码：2018MNRAS.479.2702F。doi:10.1093/mnras/sty1682。ISSN 0035-8711。
\item Opitz, Daniela；Tinney, C. G.；Faherty, Jacqueline；Sweet, Sarah；Gelino, Christopher R.；Kirkpatrick, J. Davy（2016-02-24）。“利用双子座多共轭自适应光学系统（GeMS）搜索 Y 型矮星双星”。《天体物理学杂志》（The Astrophysical Journal），819（1）：17。arXiv:1601.05508。书目代码：2016ApJ...819...17O。doi:10.3847/0004-637X/819/1/17。ISSN 1538-4357。S2CID 3208550。
\item Calissendorff, Per；De Furio, Matthew；Meyer, Michael；Albert, Loïc；Aganze, Christian；Ali-Dib, Mohamad；Bardalez Gagliuffi, Daniella C.；Baron, Frederique；Beichman, Charles A.；Burgasser, Adam J.；Cushing, Michael C.；Faherty, Jacqueline Kelly；Fontanive, Clémence；Gelino, Christopher R.；Gizis, John E.（2023-03-29）。“JWST/NIRCam 发现首个 Y+Y 型棕矮星双星：WISE J033605.05–014350.4”。《天体物理学杂志快报》（The Astrophysical Journal Letters），947（2）：L30。arXiv:2303.16923。书目代码：2023ApJ...947L..30C。doi:10.3847/2041-8213/acc86d。S2CID 257833714。
\item Bouy, Hervé。“测量超冷恒星的质量——大型地基望远镜与哈勃望远镜联合完成首个棕矮星直接质量测定”（新闻稿）。欧洲南方天文台（ESO）。检索日期：2019-12-11。
\item Bouy, Hervé；Duchêne, Gaspard；Köhler, Rainer；Brandner, Wolfgang；Bouvier, Jérôme；Martín, Eduardo L.；Ghez, Andrea Mia；Delfosse, Xavier；Forveille, Thierry；Allard, France；Baraffe, Isabelle；Basri, Gibor；Close, Laird M.；McCabe, Caer E.（2004-08-01）。“首次测定一对 L 型矮星双星的动力学质量”。《天文学与天体物理学》（Astronomy & Astrophysics），423（1）：341–352。arXiv:astro-ph/0405111。书目代码：2004A&A...423..341B。doi:10.1051/0004-6361:20040551。ISSN 0004-6361。S2CID 3149721。
\item Bedin, Luigi R.；Pourbaix, Dimitri；Apai, Dániel；Burgasser, Adam J.；Buenzli, Esther；Boffin, Henri M. J.；Libralato, Mattia（2017-09-01）。“最近棕矮星双星系统的哈勃空间望远镜天体测量——I：概述与改进的轨道”。《皇家天文学会月刊》（Monthly Notices of the Royal Astronomical Society），470（1）：1140–1155。arXiv:1706.00657。doi:10.1093/mnras/stx1177。hdl:10150/625503。ISSN 0035-8711。S2CID 119385778。
\item Luhman, Kevin L.（2004-10-10）。“宽分离棕矮星双星的首次发现”。《天体物理学杂志》（The Astrophysical Journal），614（1）：398–403。arXiv:astro-ph/0407344。书目代码：2004ApJ...614..398L。doi:10.1086/423666。ISSN 0004-637X。S2CID 11733526。
\item Reipurth, Bo；Clarke, Cathie（2003 年 6 月）。“作为被抛射恒星胚胎的棕矮星：观测视角”。IAUS，211：13–22。arXiv:astro-ph/0209005。书目代码：2003IAUS..211...13R。doi:10.1017/S0074180900210188。ISSN 1743-9221。S2CID 16822178。
\item Faherty, Jacqueline K.；Goodman, Sam；Caselden, Dan；Colin, Guillaume；Kuchner, Marc J.；Meisner, Aaron M.；Gagné, Jonathan；Schneider, Adam C.；Gonzales, Eileen C.；Bardalez Gagliuffi, Daniella C.；Logsdon, Sarah E.（2020）。“WISE2150-7520AB：通过公民科学项目 Backyard Worlds: Planet 9 发现的一套极低质量、宽分离共动棕矮星系统”。《天体物理学杂志》（The Astrophysical Journal），889（2）：176。arXiv:1911.04600。书目代码：2020ApJ...889..176F。doi:10.3847/1538-4357/ab5303。S2CID 207863267。
\item Stassun, Keivan G.；Mathieu, Robert D.；Valenti, Jeff A.（2006 年 3 月）。“在一套食双星系统中发现两颗年轻的棕矮星”。《自然》（Nature），440（7082）：311–314。书目代码：2006Natur.440..311S。doi:10.1038/nature04570。ISSN 0028-0836。PMID 16541067。S2CID 4310407。
\item Stassun, Keivan G.；Mathieu, Robert D.；Valenti, Jeff A.（2007）。“棕矮星食双星 2MASS J05352184-0546085 中温度的反常反转现象”。《天体物理学杂志》（The Astrophysical Journal），664（2）：1154–1166。arXiv:0704.3106。书目代码：2007ApJ...664.1154S。doi:10.1086/519231。S2CID 15144741。
\item Grether, Daniel；Lineweaver, Charles H.（2006-04-01）。“棕矮星荒漠有多‘干’？——量化近邻类太阳恒星周围行星、棕矮星与恒星伴星的相对数量”。《天体物理学杂志》（The Astrophysical Journal），640（2）：1051–1062。arXiv:astro-ph/0412356。书目代码：2006ApJ...640.1051G。doi:10.1086/500161。ISSN 0004-637X。S2CID 8563521。
\item Page, Emma；Pepper, Joshua；Kane, Stephen；Zhou, George；Addison, Brett；Wright, Duncan；Wittenmyer, Robert；Johnson, Marshall；Evans, Philip；Collins, Karen；Hellier, Coel；Jensen, Eric；Stassun, Keivan；Rodriguez, Joseph（2022-06-01）。“TOI-1994b：一颗凌日于次巨星前的高偏心率棕矮星”。美国天文学会会议摘要（American Astronomical Society Meeting Abstracts），54（6）：305.21。书目代码：2022AAS...24030521P。
\item “纳入系外行星档案的判据”。exoplanetarchive.ipac.caltech.edu。检索日期：2023-04-16。
\item “系外行星工作组”。w.astro.berkeley.edu。检索日期：2023-04-16。
\item Farihi, Jay；Christopher, Micol（2004 年 10 月）。“白矮星 GD 1400 可能存在一颗棕矮星伴星”。《天文学杂志》（The Astronomical Journal），128（4）：1868。arXiv:astro-ph/0407036。书目代码：2004AJ....128.1868F。doi:10.1086/423919。ISSN 1538-3881。S2CID 119530628。
\item Meisner, Aaron M.；Faherty, Jacqueline K.；Kirkpatrick, J. Davy；Schneider, Adam C.；Caselden, Dan；Gagné, Jonathan；Kuchner, Marc J.；Burgasser, Adam J.；Casewell, Sarah L.；Debes, John H.；Artigau, Étienne；Bardalez Gagliuffi, Daniella C.；Logsdon, Sarah E.；Kiman, Rocio；Allers, Katelyn（2020-08-01）。“对‘后院世界：第九行星’公民科学项目发现的极端寒冷棕矮星进行斯皮策后随观测”。《天体物理学杂志》（The Astrophysical Journal），899（2）：123。arXiv:2008.06396。书目代码：2020ApJ...899..123M。doi:10.3847/1538-4357/aba633。ISSN 0004-637X。S2CID 221135837。
\item French, Jenni R.；Casewell, Sarah L.；Dupuy, Trent J.；Debes, John H.；Manjavacas, Elena；Martin, Emily C.；Xu, Siyi（2023-03-01）。“发现一对可分辨的白矮星—棕矮星双星，投影分离很小：SDSS J222551.65+001637.7AB”。《皇家天文学会月刊》（Monthly Notices of the Royal Astronomical Society），519（4）：5008–5016。arXiv:2301.02101。书目代码：2023MNRAS.519.5008F。doi:10.1093/mnras/stac3807。ISSN 0035-8711。
\item Leggett, S. K.；Tremblin, P.；Esplin, T. L.；Luhman, K. L.；Morley, Caroline V.（2017-06-01）。“Y 型棕矮星：基于新的天体测量、同质化测光与近红外光谱对质量和年龄的估计”。《天体物理学杂志》（The Astrophysical Journal），842（2）：118。arXiv:1704.03573。书目代码：2017ApJ...842..118L。doi:10.3847/1538-4357/aa6fb5。ISSN 0004-637X。S2CID 119249195。
\item Maxted, Pierre；Napiwotzki, Ralf；Dobbie, Paul；Burleigh, Matt。“一位亚恒星的约拿——棕矮星在被吞噬后幸存”（新闻稿）。欧洲南方天文台（European Southern Observatory）。检索日期：2019-12-11。
\item Casewell, Sarah L.；Braker, Ian P.；Parsons, Steven G.；Hermes, James J.；Burleigh, Matthew R.；Belardi, Claudia；Chaushev, Alexander；Finch, Nicolle L.；Roy, Mervyn；Littlefair, Stuart P.；Goad, Mike；Dennihy, Erik（2018 年 1 月 31 日）。“首个周期短于 70 分钟的非相互作用白矮星—棕矮星系统：EPIC212235321”。《皇家天文学会月刊》（Monthly Notices of the Royal Astronomical Society），476（1）：1405–1411。arXiv:1801.07773。书目代码：2018MNRAS.476.1405C。doi:10.1093/mnras/sty245。ISSN 0035-8711。S2CID 55776991。
\item Longstaff, Emma S.；Casewell, Sarah L.；Wynn, Graham A.；Maxted, Pierre F. L.；Helling, Christiane（2017-10-21）。“受辐照棕矮星 WD0137−349B 大气中的发射线”。《皇家天文学会月刊》（Monthly Notices of the Royal Astronomical Society），471（2）：1728–1736。arXiv:1707.05793。书目代码：2017MNRAS.471.1728L。doi:10.1093/mnras/stx1786。ISSN 0035-8711。S2CID 29792989。
\item Grether, Daniel；Lineweaver, Charles H.（2006 年 4 月）。“棕矮星荒漠有多‘干’？——量化近邻类太阳恒星周围行星、棕矮星与恒星伴星的相对数量”。《天体物理学杂志》（The Astrophysical Journal），640（2）：1051–1062。arXiv:astro-ph/0412356。书目代码：2006ApJ...640.1051G。doi:10.1086/500161。ISSN 0004-637X。
\item Rappaport, Saul A.；Vanderburg, Andrew；Nelson, Lorne；Gary, Bruce L.；Kaye, Thomas G.；Kalomeni, Belinda；Howell, Steve B.；Thorstensen, John R.；Lachapelle, François-René；Lundy, Matthew；St-Antoine, Jonathan（2017-10-11）。“WD 1202-024：周期最短的灾变前变星”。《皇家天文学会月刊》（Monthly Notices of the Royal Astronomical Society），471（1）：948–961。arXiv:1705.05863。书目代码：2017MNRAS.471..948R。doi:10.1093/mnras/stx1611。ISSN 0035-8711。S2CID 119349942。
\item Neustroev, Vitaly V.；Mäntynen, Iikka（2023-08-01）。“周期回弹（period-bouncer）候选体 BW Sculptoris：一颗棕矮星供体、一个光学薄的吸积盘，以及具有复杂流撞击区域的系统”。《皇家天文学会月刊》（Monthly Notices of the Royal Astronomical Society），523（4）：6114–6137。arXiv:2212.03264。书目代码：2023MNRAS.523.6114N。doi:10.1093/mnras/stad1730。ISSN 0035-8711。
\item Kato, Taichi（2015-12-01）。“WZ Sge 型矮新星”。《日本天文学会出版物》（Publications of the Astronomical Society of Japan），67（6）：108。arXiv:1507.07659。书目代码：2015PASJ...67..108K。doi:10.1093/pasj/psv077。ISSN 0004-6264。
\item Lira, Nicolás；Blue, Charles E.；Turner, Calum；Hiramatsu, Masaaki。“什么时候新星不是‘新星’？——当白矮星与棕矮星相撞时”。ALMA 天文台。原文档于 2019-10-22 存档。检索日期：2019-11-12。
\item Eyres, Stewart P. S.；Evans, Aneurin；Zijlstra, Albert；Avison, Adam；Gehrz, Robert D.；Hajduk, Marcin；Starrfield, Sumner；Mohamed, Shazrene；Woodward, Charles E.；Wagner, R. Mark（2018-12-21）。“ALMA 揭示白矮星—棕矮星并合事件在 CK Vulpeculae 中的余波”。《皇家天文学会月刊》（Monthly Notices of the Royal Astronomical Society），481（4）：4931–4939。arXiv:1809.05849。书目代码：2018MNRAS.481.4931E。doi:10.1093/mnras/sty2554。ISSN 0035-8711。S2CID 119462149。
\item Palau, Aina；Huelamo, Nuria；Barrado, David；Dunham, Michael M.；Lee, Chang Won（2024-10-01）。“近邻分子云中前棕矮星与原棕矮星的观测：为进一步约束棕矮星形成理论铺路”。《新天文学评论》（New Astronomy Reviews），99：101711。arXiv:2410.07769。书目代码：2024NewAR..9901711P。doi:10.1016/j.newar.2024.101711。检索日期：2024-10-25。
\item Riaz, Basmah；Machida, Masahiro N.；Stamatellos, Dimitris（2019 年 7 月）。“ALMA 揭示一颗原棕矮星中的伪盘（pseudo-disc）”。《皇家天文学会月刊》（Monthly Notices of the Royal Astronomical Society），486（3）：4114–4129。arXiv:1904.06418。书目代码：2019MNRAS.486.4114R。doi:10.1093/mnras/stz1032。ISSN 0035-8711。S2CID 119286540。
\item Riaz, Basmah；Najita, Joan。“超出其‘体量’的表现：一颗棕矮星发射出秒差距尺度喷流”。美国国家光学天文台（National Optical Astronomy Observatory）。原文档于 2020-02-18 存档。检索日期：2020-02-18。
\item Riaz, Basmah；Briceño, Cesar；Whelan, Emma T.；Heathcote, Stephen（2017 年 7 月）。“由原棕矮星驱动的首个大尺度 Herbig–Haro 喷流”。《天体物理学杂志》（Astrophysical Journal），844（1）：47。arXiv:1705.01170。书目代码：2017ApJ...844...47R。doi:10.3847/1538-4357/aa70e8。ISSN 0004-637X。S2CID 119080074。
\item Apai, Dániel；Pascucci, Ilaria；Bouwman, Jeroen；Natta, Antonella；Henning, Thomas；Dullemond, Cornelis P.（2005）。“棕矮星盘中行星形成的起始”。《科学》（Science），310（5749）：834–836。arXiv:astro-ph/0511420。书目代码：2005Sci...310..834A。doi:10.1126/science.1118042。PMID 16239438。S2CID 5181947。
\item Burrows, Adam；Hubbard, William B.；Lunine, Jonathan I.；Liebert, James（2011）。“棕矮星周围行星的潮汐演化”。《天文学与天体物理学》（Astronomy & Astrophysics），535：A94。arXiv:1109.2906。书目代码：2011A&A...535A..94B。doi:10.1051/0004-6361/201117734。S2CID 118532416。
\item Jewitt, David C.，《Pan-STARRS 科学概览》。原文档于 2015-10-16 存档于 Wayback Machine。
\item Schutte, Maria（2020-08-12）。“我们的新论文：发现一颗邻近的年轻棕矮星盘！”DiskDetective.org。检索日期：2023-09-23。
\item Schutte, Maria C.；Lawson, Kellen D.；Wisniewski, John P.；Kuchner, Marc J.；Silverberg, Steven M.；Faherty, Jacqueline K.；Gagliuffi, Daniella C. Bardalez；Kiman, Rocio；Gagné, Jonathan；Meisner, Aaron；Schneider, Adam C.；Bans, Alissa S.；Debes, John H.；Kovacevic, Natalie；Bosch, Milton K. D.；Luca, Hugo A. Durantini；Holden, Jonathan；Hyogo, Michiharu（2020-08-04）。“发现一颗邻近的年轻棕矮星盘”。《天文学杂志》（The Astronomical Journal），160（4）：156。arXiv:2007.15735v2。书目代码：2020AJ....160..156S。doi:10.3847/1538-3881/abaccd。ISSN 1538-3881。S2CID 220920317。
\item Rilinger, Anneliese M.；Espaillat, Catherine C.（2021-11-01）。“棕矮星周围原行星盘的盘质量与尘埃演化”。《天体物理学杂志》（The Astrophysical Journal），921（2）：182。arXiv:2106.05247。书目代码：2021ApJ...921..182R。doi:10.3847/1538-4357/ac09e5。ISSN 0004-637X。S2CID 235377000。
\item Luhman, K. L.；Alves de Oliveira, C.；Baraffe, I.；Chabrier, G.；Manjavacas, E.；Parker, R. J.；Tremblin, P.（2024 年 10 月 13 日）。“JWST/NIRSpec 对猎户座星云星团棕矮星的观测”。《天体物理学杂志》（The Astrophysical Journal），975（2）：162。arXiv:2410.10000。书目代码：2024ApJ...975..162L。doi:10.3847/1538-4357/ad7b19。
\item Ricci, L.；Testi, L.；Natta, A.；Scholz, A.；de Gregorio-Monsalvo, I.；Isella, A.（2014-08-01）。“ALMA 观测的棕矮星盘”。《天体物理学杂志》（The Astrophysical Journal），791（1）：20。arXiv:1406.0635。书目代码：2014ApJ...791...20R。doi:10.1088/0004-637X/791/1/20。hdl:10023/8607。ISSN 0004-637X。S2CID 13180928。
\item Boucher, Anne；Lafrenière, David；Gagné, Jonathan；Malo, Lison；Faherty, Jacqueline K.；Doyon, René；Chen, Christine H.（2016-11-01）。“BANYAN VIII：带有候选环星盘的新低质量恒星与棕矮星”。《天体物理学杂志》（The Astrophysical Journal），832（1）：50。arXiv:1608.08259。书目代码：2016ApJ...832...50B。doi:10.3847/0004-637X/832/1/50。ISSN 0004-637X。S2CID 119017727。
\item Silverberg, Steven M.；Wisniewski, John P.；Kuchner, Marc J.；Lawson, Kellen D.；Bans, Alissa S.；Debes, John H.；Biggs, Joseph R.；Bosch, Milton K. D.；Doll, Katharina；Luca, Hugo A. Durantini；Enachioaie, Alexandru；Hamilton, Joshua；Holden, Jonathan；Hyogo, Michiharu（2020-02-01）。“彼得·潘盘：年轻 M 型恒星周围的长寿吸积盘”。《天体物理学杂志》（The Astrophysical Journal），890（2）：106。arXiv:2001.05030。书目代码：2020ApJ...890..106S。doi:10.3847/1538-4357/ab68e6。ISSN 0004-637X。S2CID 210718358。
\item Luhman, Kevin L.；Adame, Lucía；d’Alessio, Paola；Calvet, Nuria；Hartmann, Lee；Megeath, S. T.；Fazio, G. G.（2005）。“发现一颗具有环星盘的行星质量棕矮星”。《天体物理学杂志》（The Astrophysical Journal），635（1）：L93–L96。arXiv:astro-ph/0511807。书目代码：2005ApJ...635L..93L。doi:10.1086/498868。S2CID 11685964。
\item Ricci, Luca；Testi, Leonardo；Pierce-Price, Douglas；Stoke, John。“即便是棕矮星也可能孕育岩质行星”（新闻稿）。欧洲南方天文台。原文档于 2012 年 12 月 3 日存档。检索日期：2012 年 12 月 3 日。
\item Lecavelier des Etangs, A.；Lissauer, Jack J.（2022-06-01）。“IAU 对系外行星的工作定义”。《新天文学评论》（New Astronomy Reviews），94：101641。arXiv:2203.09520。书目代码：2022NewAR..9401641L。doi:10.1016/j.newar.2022.101641。ISSN 1387-6473。S2CID 247065421。
\item Joergens, Viki；Müller, André（2007）。“围绕棕矮星候选体 Cha Hα 8 的 16–20 $M_{\mathrm{Jup}}$ 径向速度伴星”。《天体物理学杂志》（The Astrophysical Journal），666（2）：L113–L116。arXiv:0707.3744。书目代码：2007ApJ...666L.113J。doi:10.1086/521825。S2CID 119140521。
\item Joergens, Viki；Müller, André；Reffert, Sabine（2010）。“年轻双星棕矮星候选体 Cha Hα 8 的改进径向速度轨道解”。《天文学与天体物理学》（Astronomy and Astrophysics），521（A24）：A24。arXiv:1006.2383。书目代码：2010A&A...521A..24J。doi:10.1051/0004-6361/201014853。S2CID 54989533。
\item Bennet, David P.；Bond, Ian A.；Udalski, Andrzej；Sumi, Takahiro；Abe, Fumio；Fukui, Akihiko；Furusawa, Kei；Hearnshaw, John B.；Holderness, Sarah；Itow, Yoshitaka；Kamiya, Koki；Korpela, Aarno V.；Kilmartin, Pamela M.；Lin, Wei；Ling, Cho Hong；Masuda, Kimiaki；Matsubara, Yutaka；Miyake, Noriyuki；Muraki, Yasushi；Nagaya, Maiko；Okumura, Teppei；Ohnishi, Kouji；Perrott, Yvette C.；Rattenbury, Nicholas J.；Sako, Takashi；Saito, Toshiharu；Sato, S.；Skuljan, Ljiljana；Sullivan, Denis J.；Sweatman, Winston L.；Tristram, Paul J.；Yock, Philip C. M.；Kubiak, Marcin；Szymański, Michał K.；Pietrzyński, Grzegorz；Soszyński, Igor；Szewczyk, O.；Wyrzykowski, Łukasz；Ulaczyk, Krzysztof；Batista, Virginie；Beaulieu, Jean-Philippe；Brillant, Stéphane；Cassan, Arnaud；Fouqué, Pascal；Kervella, Pierre；Kubas, Daniel；Marquette, Jean-Baptiste（2008 年 5 月 30 日）。“微引力透镜事件 MOA-2007-BLG-192 中的一颗低质量行星：其宿主可能为亚恒星质量”。《天体物理学杂志》（The Astrophysical Journal），684（1）：663–683。arXiv:0806.0025。书目代码：2008ApJ...684..663B。doi:10.1086/589940。S2CID 14467194。
\item Burrows, Adam；Hubbard, William B.；Lunine, Jonathan I.；Liebert, James（2013）。“极低质量恒星与棕矮星周围盘的原子与分子成分”。《天体物理学杂志》（The Astrophysical Journal），779（2）：178。arXiv:1311.1228。书目代码：2013ApJ...779..178P。doi:10.1088/0004-637X/779/2/178。S2CID 119001471。
\item He, Matthias Y.；Triaud, Amaury H. M. J.；Gillon, Michaël（2017 年 1 月）。“棕矮星周围短周期行星出现率的首批上限”。《皇家天文学会月刊》（Monthly Notices of the Royal Astronomical Society），464（3）：2687–2697。arXiv:1609.05053。书目代码：2017MNRAS.464.2687H。doi:10.1093/mnras/stw2391。S2CID 53692008。
\item Limbach, Mary Anne；Vos, Johanna M.；Winn, Joshua N.；Heller, René；Mason, Jeffrey C.；Schneider, Adam C.；Dai, Fei（2021-09-01）。“关于探测凌日孤立行星质量天体的系外卫星”。《天体物理学杂志》（The Astrophysical Journal），918（2）：L25。arXiv:2108.08323。书目代码：2021ApJ...918L..25L。doi:10.3847/2041-8213/ac1e2d。ISSN 0004-637X。
\item Vos, Johanna M.；Allers, Katelyn N.；Biller, Beth A.（2017-06-01）。“棕矮星的观测几何会影响其观测到的颜色与变幅”。《天体物理学杂志》（The Astrophysical Journal），842（2）：78。arXiv:1705.06045。书目代码：2017ApJ...842...78V。doi:10.3847/1538-4357/aa73cf。ISSN 0004-637X。
\item Baycroft, Thomas A.；Sairam, Lalitha；Triaud, Amaury H. M. J.；Correia, Alexandre C. M.（2025-04-16）。“证据表明：一颗极向环双（polar circumbinary）系外行星绕一对食双棕矮星运行”。《科学·进展》（Science Advances），11（16）：eadu0627。arXiv:2504.12209。书目代码：2025SciA...11..627B。doi:10.1126/sciadv.adu0627。PMC 12002110。PMID 40238865。
\item “‘大惊喜’：天文学家发现行星以垂直轨道绕一对恒星运行”。[www.eso.org。检索日期：2025-04-16。](http://www.eso.org。检索日期：2025-04-16。)
\item Barnes, Rory；Heller, René（2011）。“白矮星与棕矮星周围的宜居行星：主星冷却带来的风险”。《天体生物学》（Astrobiology），13（3）：279–291。arXiv:1211.6467。书目代码：2013AsBio..13..279B。doi:10.1089/ast.2012.0867。PMC 3612282。PMID 23537137。
\item Morrison, David（2011 年 8 月 2 日）。“如今科学家不再认为类似‘复仇女神（Nemesis）’的天体可能存在”。NASA Ask An Astrobiologist。原文档于 2012 年 12 月 13 日存档。检索日期：2011-10-22。
\item Burrows, Adam；Hubbard, William B.；Lunine, Jonathan I.；Liebert, James（2004）。“用哈勃空间望远镜在 L 型矮星双星系统 DENIS-P J020529.0-115925 中发现可能的第三个成员”。《天文学杂志》（The Astronomical Journal），129（1）：511–517。arXiv:astro-ph/0410226。书目代码：2005AJ....129..511B。doi:10.1086/426559。S2CID 119336794。
\item Basri, Gibor；Martín, Eduardo L.（1999）。“PPl 15：首个棕矮星光谱双星”。《天文学杂志》（The Astronomical Journal），118（5）：2460–2465。arXiv:astro-ph/9908015。书目代码：1999AJ....118.2460B。doi:10.1086/301079。S2CID 17662168。
\item Comerón, F.；Neuhäuser, R.；Kaas, A. A.（2000-07-01）。“探测蝘蜒座 I 恒星形成区的棕矮星族群”。《天文学与天体物理学》（Astronomy and Astrophysics），359：269–288。书目代码：2000A&A...359..269C。ISSN 0004-6361。
\item Scholz, Ralf-Dieter；McCaughrean, Mark（2003-01-13）。“eso0303——发现迄今已知最近的棕矮星”（新闻稿）。欧洲南方天文台（European Southern Observatory）。原文档于 2008 年 10 月 13 日存档。检索日期：2013-03-16。
\item Burgasser, Adam J.；Kirkpatrick, J. Davy；Burrows, Adam；Liebert, James；Reid, I. Neill；Gizis, John E.；McGovern, Mark R.；Prato, Lisa；McLean, Ian S.（2003）。“首个亚恒星次矮星？——发现一颗具有晕族运动学特征的贫金属 L 型矮星”。《天体物理学杂志》（The Astrophysical Journal），592（2）：1186–1192。arXiv:astro-ph/0304174。书目代码：2003ApJ...592.1186B。doi:10.1086/375813。S2CID 11895472。
\item Whelan, Emma T.；Ray, Thomas P.；Bacciotti, Francesca；Natta, Antonella；Testi, Leonardo；Randich, Sofia（2005 年 6 月）。“来自一颗棕矮星的可分辨物质外流”。《自然》（Nature），435（7042）：652–654。arXiv:astro-ph/0506485。书目代码：2005Natur.435..652W。doi:10.1038/nature03598。ISSN 0028-0836。PMID 15931217。S2CID 4415442。
\item Maxted, Pierre F. L.；Napiwotzki, Ralf；Dobbie, Paul D.；Burleigh, Matthew R.（2006）。“棕矮星在被红巨星吞没后幸存”。《自然》（Nature，投稿手稿），442（7102）：543–545。arXiv:astro-ph/0608054。书目代码：2006Natur.442..543M。doi:10.1038/nature04987。hdl:2299/1227。PMID 16885979。S2CID 4368344。原文档于 2021-04-27 存档。
\item Wolszczan, Alexander；Route, Matthew（2014）。“对超冷矮星 TVLM 513-46546 的周期性射电与光学亮度变化的计时分析”。《天体物理学杂志》（The Astrophysical Journal），788（1）：23。arXiv:1404.4682。书目代码：2014ApJ...788...23W。doi:10.1088/0004-637X/788/1/23。S2CID 119114679。
\item Vallenari, A. 等（Gaia 合作组）（2023）。“Gaia 数据发布 3：内容与巡天性质概述”。《天文学与天体物理学》（Astronomy and Astrophysics），674：A1。arXiv:2208.00211。书目代码：2023A&A...674A...1G。doi:10.1051/0004-6361/202243940。S2CID 244398875。该源在 VizieR 的 Gaia DR3 记录。
\item Roman, Nancy G.（1987）。“由一个位置识别所属星座”。《太平洋天文学会出版物》（Publications of the Astronomical Society of the Pacific），99（617）：695。书目代码：1987PASP...99..695R。doi:10.1086/132034。该天体在 VizieR 的星座记录。
\item Redai, Jéa Adams；Chandra, Vedant；Yee, Samuel W.；DiTomasso, Victoria；Andrews, Sean；Öberg, Karin；Woody, Rebecca；Latham, David W.；Bieryla, Allyson（2025-12-05）。“一颗古老棕矮星凌日一颗贫金属厚盘恒星”。arXiv:2512.06069 [astro-ph.EP]。
\item Levine, Joanna L.；Steinhauer, Aaron；Elston, Richard J.；Lada, Elizabeth A.（2006-08-01）。“NGC 2024 中的低质量恒星与棕矮星：对亚恒星质量函数的约束”。《天体物理学杂志》（The Astrophysical Journal），646（2）：1215–1229。arXiv:astro-ph/0604315。书目代码：2006ApJ...646.1215L。doi:10.1086/504964。ISSN 0004-637X。S2CID 118955538。表 3：FLMN_J0541328-0151271。
\item Zhang, ZengHua；Homeier, Derek；Pinfield, David J.；Lodieu, Nicolas；Jones, Hugh R. A.；Pavlenko, Yakiv V.（2017-06-11）。“原初极低质量恒星与棕矮星 II：金属丰度最低的亚恒星天体”。《皇家天文学会月刊》（Monthly Notices of the Royal Astronomical Society），468（1）：261。arXiv:1702.02001。书目代码：2017MNRAS.468..261Z。doi:10.1093/mnras/stx350。S2CID 54847595。
\item Luhman, Kevin（2009）。“发现一对在孤立环境中诞生的宽分离棕矮星双星”。《天体物理学杂志》（The Astrophysical Journal），691（2）：1265。arXiv:0902.0425。书目代码：2009ApJ...691.1265L。doi:10.1088/0004-637X/691/2/1265。
\item “根据给定的天球坐标查找所属星座”。djm.cc。检索日期：2025-01-16。
\item Burdge, Kevin B.；Marsh, Thomas R.；Fuller, Jim；Bellm, Eric C.；Caiazzo, Ilaria；Chakrabarty, Deepto；Coughlin, Michael W.；De, Kishalay；Dhillon, V. S.；Graham, Matthew J.；Rodríguez-Gil, Pablo；Jaodand, Amruta D.；Kaplan, David L.；Kara, Erin；Kong, Albert K. H.（2022 年 5 月）。“宽层级三体系统中的一颗 62 分钟轨道周期‘黑寡妇’双星”。《自然》（Nature），605（7908）：41–45。arXiv:2205.02278。书目代码：2022Natur.605...41B。doi:10.1038/s41586-022-04551-1。ISSN 1476-4687。PMID 35508781。
\item Tannock, Megan E.；Metchev, Stanimir；Heinze, Aren；Miles-Páez, Paulo A.；Gagné, Jonathan；Burgasser, Adam J.；Marley, Mark S.；Apai, Dániel；Suárez, Genaro；Plavchan, Peter（2021 年 3 月）。“其他世界的天气 V：自转最快的三颗超冷矮星”。《天文学杂志》（The Astronomical Journal），161（5）：224。arXiv:2103.01990。书目代码：2021AJ....161..224T。doi:10.3847/1538-3881/abeb67。S2CID 232105126。
\item Zeidler, Peter；Sabbi, Elena；Nota, Antonella；Manjavacas, Elena；Jones, Olivia C.；Pacifici, Camilla（2024）。“在小麦哲伦云中发现亚太阳金属丰度的棕矮星候选体”。《天体物理学杂志》（The Astrophysical Journal），975（1）：18。书目代码：2024ApJ...975...18Z。doi:10.3847/1538-4357/ad779e。
\item Sanghi, Aniket；Liu, Michael C.；Best, William M. J.；Dupuy, Trent J.；Siverd, Robert J.；Zhang, Zhoujian；Hurt, Spencer A.；Magnier, Eugene A.；Aller, Kimberly M.；Deacon, Niall R.（2023）。“超冷天体基本物理性质表”。Zenodo。doi:10.5281/zenodo.10086810。
\item Rose, Kovi；Pritchard, Joshua；Murphy, Tara；Caleb, Manisha；Dobie, Dougal；Driessen, Laura；Duchesne, Stefan；Kaplan, David；Lenc, Emil；Wang, Ziteng（2023）。“T8 型矮星 WISE J062309.94−045624.6 的周期性射电辐射”。《天体物理学杂志快报》（The Astrophysical Journal Letters），951（2）：L43。arXiv:2306.15219。书目代码：2023ApJ...951L..43R。doi:10.3847/2041-8213/ace188。S2CID 259262475。
\item Astrobites（2020 年 6 月 24 日）。“TESS 发现的凌日棕矮星 2”。AAS Nova。检索日期：2013-03-16。
\end{enumerate}
\subsection{外部链接}
\begin{itemize}
\item HubbleSite 新闻中心——棕矮星上的天气模式
\item Allard, France；Homeier, Derek（2007）。“棕矮星”。Scholarpedia，2（12）：4475。书目代码：2007SchpJ...2.4475A。doi:10.4249/scholarpedia.4475。
\end{itemize}
\subsubsection{历史}
\begin{itemize}
\item Kumar, Shiv S.；*Low-Luminosity Stars*。Gordon and Breach，伦敦，1969——一篇关于棕矮星的早期综述论文
《《哥伦比亚百科全书》：“棕矮星（Brown Dwarfs）”
\end{itemize}
\subsubsection{详细信息}
\begin{itemize}
\item L 型与 T 型矮星的最新列表
一种“地质学”意义下对棕矮星的定义，并与恒星与行星进行对比（经由伯克利）
\item I. Neill Reid 在空间望远镜科学研究所（Space Telescope Science Institute）的页面：
\item 关于 M 型、L 型与 T 型矮星的光谱分析
\item 低温矮星的温度与质量特征
\item 首次观测到来自棕矮星的 X 射线，Spaceref.com，2000
\item Montes, David；“棕矮星与超冷矮星（晚期 M、L、T）”，马德里康普顿斯大学（UCM）
\item 《狂野天气：失败恒星上的铁雨》——科学家正在研究棕矮星中令人震惊的天气模式，Space.com，2006
\item NASA “棕矮星侦探”（Brown dwarf detectives）原文档于 2014-10-17 存档于 Wayback Machine——以更简化的方式提供更详尽的信息
\item “棕矮星（Brown Dwarfs）”——提供棕矮星一般信息的网站（包含许多细节丰富、色彩鲜明的艺术家想象图）
\end{itemize}
\subsubsection{恒星}
\begin{itemize}
\item Cha Halpha 1 的数据与历史
\item “已观测棕矮星的统计”（并非全部已被确认），1998
\item Luhman, Kevin L.；Adame, Lucía；d’Alessio, Paola；Calvet, Nuria；Hartmann, Lee；Megeath, S. Thomas；Fazio, Giovanni G.（2005）。“发现一颗具有环星盘的行星质量棕矮星”。《天体物理学杂志》（The Astrophysical Journal），635（1）：L93–L96。arXiv:astro-ph/0511807。书目代码：2005ApJ...635L..93L。doi:10.1086/498868。S2CID 11685964。
\item Michaud, Peter；Heyer, Inge；Leggett, Sandy K.；Adamson, Andy；“这一发现缩小了行星与棕矮星之间的鸿沟”，Gemini 与联合天文中心（Joint Astronomy Centre），2007
\item Deacon, N. R.；Hambly, N. C.（2006）。“用 UKIRT 红外深空巡天（UKIDSS）探测超冷矮星的可能性”。《皇家天文学会月刊》（Monthly Notices of the Royal Astronomical Society），371（4）：1722–1730。arXiv:astro-ph/0607305。书目代码：2006MNRAS.371.1722D。doi:10.1111/j.1365-2966.2006.10795.x。
\end{itemize}